% From matplotlib and seaborn to simplydata: written by hand, read by docs/make_manual.py (% From matplotlib and seaborn to simplydata: written by hand, read by docs/make_manual.py (% From matplotlib and seaborn to simplydata: written by hand, read by docs/make_manual.py (% From matplotlib and seaborn to simplydata: written by hand, read by docs/make_manual.py (\input{from_matplotlib}).
% The topics follow the order of the matplotlib Quick start guide and User guide, then the seaborn tutorial, used as a
% checklist only: the text and the examples are our own, and each section links to the official page. Every example in
% a py box is run by tests/test_from_matplotlib_chapter.py (with "import simplydata as sdata" in front, and matplotlib's
% Agg backend), so it cannot go out of date. A line "% needs: plotly" just before a box says the box needs that package.

\unchapter{From matplotlib and seaborn to simplydata}{from-matplotlib}

This chapter is for you if you already know matplotlib or seaborn, or if you want to learn them after
simplydata. It walks through the topics of the official guides, in their order: the
\href{https://matplotlib.org/stable/users/explain/quick_start.html}{matplotlib Quick start guide} and
\href{https://matplotlib.org/stable/users/index.html}{User guide}, then the
\href{https://seaborn.pydata.org/tutorial.html}{seaborn tutorial}. For each topic it says what the job
is, which simplydata words do it, and a short example on the practice tables. Some jobs are left to
matplotlib and seaborn on purpose. For those, the chapter shows how to hand the chart over and keep
going in matplotlib.

Every example starts with \code{import simplydata as sdata}. It is written once here, not in each
box. The examples use the practice tables of \code{sdata.get\_example}, so you can run them as they
are.

\unsection{Five things to know first}{from-matplotlib:first}

\begin{enumerate}
  \item \textbf{A chart is a list of steps.} In matplotlib each call draws at once. In simplydata,
    \code{chart.draw\_dots(...)} and \code{chart.label\_x(...)} only write down a step. The chart is
    drawn when you save it, show it, or look at it in Jupyter. So the order of the steps rarely
    matters, and \code{chart.undo()} takes the last step back.
  \item \textbf{There is no figure and no axes to hold.} simplydata makes the matplotlib Figure and
    its Axes for you each time it draws. You never write \code{fig, ax = plt.subplots()}.
  \item \textbf{Columns by name.} You give the table once, to \code{sdata.chart(table)}, and then the
    column names, as seaborn does with \code{data=}. A word works on a pandas or a Polars table.
  \item \textbf{One word for one job.} matplotlib often has several ways to do the same thing
    (\code{plt.title}, \code{ax.set\_title}, \code{fig.suptitle}). simplydata has one word,
    \code{write\_title\_on\_top}, and \code{sdata.help("ax.set\_title")} names it.
  \item \textbf{The door is always open.} \code{fig = chart.as\_matplotlib()} gives the chart as a
    plain matplotlib Figure. Its charts are the Axes in \code{fig.axes}, so \code{ax = fig.axes[0]}.
    Then any matplotlib code works. It is a copy: changing \code{fig} does not change the chart.
    Save it with \code{fig.savefig(...)}.
\end{enumerate}

\begin{lstlisting}[style=py]
marks = sdata.get_example("marks")
chart = sdata.chart(marks)
chart.draw_dots(x="hours", y="marks")
fig = chart.as_matplotlib()              # a matplotlib Figure: use any matplotlib code here
ax = fig.axes[0]
ax.set_facecolor("whitesmoke")
fig.savefig("handed_over.png", dpi=150)
\end{lstlisting}

Do you know a matplotlib or seaborn name, but not the simplydata word? Ask:
\code{sdata.help("ax.set\_xlabel")}, \code{sdata.help("plt.hist")} or
\code{sdata.help("sns.lineplot")}. The name can have \code{plt.}, \code{ax.}, \code{fig.} or
\code{sns.} in front. The full list is in the section ``Coming from another tool?'' of the simplydata
dictionary
(\href{https://github.com/bioforkgithub/simplydata/blob/main/docs/manual/simplydataDictionary.pdf}%
{simplydataDictionary.pdf}).

\unsection{A first chart}{from-matplotlib:first-chart}

\textbf{The job.} Draw a line from two columns, give it a title and axis names, and save it as a
picture. In the guide: \href{https://matplotlib.org/stable/users/explain/quick_start.html}{Quick
start guide}, the simple example at its start.

\textbf{In simplydata.} \code{draw\_lines} is \code{ax.plot}; \code{write\_title\_on\_top},
\code{label\_x} and \code{label\_y} are \code{set\_title}, \code{set\_xlabel} and
\code{set\_ylabel}; \code{save} is \code{fig.savefig}. Without \code{label\_x}, the column name is the
axis name.

\begin{lstlisting}[style=py]
weather = sdata.get_example("weather")
chart = sdata.chart(weather)
chart.draw_lines(x="month", y="temperature", colour="city")
chart.write_title_on_top("Temperature through the year")
chart.label_y("Temperature (degrees C)")
chart.save("first_chart.png")
\end{lstlisting}

\unsection{The parts of a figure}{from-matplotlib:parts}

\textbf{The job.} Know the parts matplotlib draws: the Figure (the whole picture), the Axes (one
chart inside it), each Axis (the x and y scales) and the Artists (every line, dot, bar and piece of
text). In the guide: \href{https://matplotlib.org/stable/users/explain/quick_start.html}{Quick start
guide}, ``Parts of a Figure'', and
\href{https://matplotlib.org/stable/users/explain/artists/index.html}{Artists}.

\textbf{In simplydata.} A simplydata chart is one Figure with one Axes, or several Axes after
\code{draw\_one\_chart\_for\_each\_group\_in}. You do not touch the parts. You can read them back as
tables: \code{get\_numbers\_drawn} reads the lines, dots and bars, \code{get\_text\_drawn} the text,
and \code{get\_scales\_drawn} each axis. \code{as\_matplotlib} gives the parts themselves.

\begin{lstlisting}[style=py]
marks = sdata.get_example("marks")
chart = sdata.chart(marks)
chart.draw_bars(x="class", y="marks")
chart.get_numbers_drawn()                # one row for each bar: where it is and how high
chart.get_scales_drawn()                 # the range and kind of each axis
fig = chart.as_matplotlib()
print(len(fig.axes), "Axes in the Figure")
\end{lstlisting}

\unsection{The data a chart takes}{from-matplotlib:data}

\textbf{The job.} Give the numbers to draw. matplotlib takes lists or NumPy arrays, or column names
with \code{data=}. In the guide: \href{https://matplotlib.org/stable/users/explain/quick_start.html}%
{Quick start guide}, ``Types of inputs to plotting functions''.

\textbf{In simplydata.} A table and column names is the usual way: \code{sdata.chart(table)} reads a
table, a file name or a pandas or Polars DataFrame. Without a table, give lists to the step itself.
Numbers made with \code{sdata.number\_list} start their own chart.

\begin{lstlisting}[style=py]
chart = sdata.chart()                    # no table: lists go to the step
chart.draw_lines(x=[1, 2, 3, 4], y=[10, 12, 9, 14])
chart.save("from_lists.png")

numbers = sdata.number_list([4, 8, 15, 16, 23, 42])
chart = numbers.start_chart()
chart.draw_histogram(x="value")
chart.save("from_numbers.png")
\end{lstlisting}

\unsection{Two ways of writing, one in simplydata}{from-matplotlib:styles}

\textbf{The job.} matplotlib can be written in two styles: the explicit one (\code{fig, ax =
plt.subplots()}, then \code{ax.plot(...)}) and the implicit one with \code{pyplot}
(\code{plt.plot(...)}, which draws on the current Axes). In the guide:
\href{https://matplotlib.org/stable/users/explain/quick_start.html}{Quick start guide}, ``Coding
styles''.

\textbf{In simplydata.} There is one style. The chart is the object you name, as in the explicit style,
and nothing is ``current''. So two charts never draw on each other. To make many charts the same way,
write a small function that gives back a chart, as the guide suggests for matplotlib.

\begin{lstlisting}[style=py]
def marks_chart(subject):
    marks = sdata.get_example("marks")
    marks.keep_rows_where("subject", is_equal_to=subject)
    chart = sdata.chart(marks)
    chart.draw_dots(x="hours", y="marks")
    chart.write_title_on_top(subject)
    return chart


marks_chart("Maths").save("maths.png")
marks_chart("Science").save("science.png")
\end{lstlisting}

\unsection{Colours, lines and dots}{from-matplotlib:styling}

\textbf{The job.} Choose the colour, width and dash of a line, and the size, shape and transparency of
dots. In the guide: \href{https://matplotlib.org/stable/users/explain/quick_start.html}{Quick start
guide}, ``Styling Artists'', and
\href{https://matplotlib.org/stable/users/explain/colors/colors.html}{Specifying colors}.

\textbf{In simplydata.} Options with plain names on the draw words: \code{thickness=} is
\code{linewidth}, \code{dashed=True} and \code{dotted=True} are \code{linestyle}, \code{size=} is the
dot size (a column, for sizes that change), \code{transparency=} is one minus \code{alpha}.
\code{use\_colour} sets one colour, \code{use\_colours} one colour for each group, and
\code{use\_shapes} a shape for each group. The groups come from \code{colour="column"}, as seaborn's
\code{hue=}.

\begin{lstlisting}[style=py]
weather = sdata.get_example("weather")
chart = sdata.chart(weather)
chart.draw_lines(x="month", y="temperature", colour="city", thickness=3, dashed="city")
chart.use_colours("black", "darkorange")
chart.save("styled_lines.png")

marks = sdata.get_example("marks")
chart = sdata.chart(marks)
chart.draw_dots(x="hours", y="marks", colour="class", size="height", transparency=0.3)
chart.use_shapes("circle", "square", "triangle")
chart.save("styled_dots.png")
\end{lstlisting}

\unsection{Titles, labels, notes and legends}{from-matplotlib:labels}

\textbf{The job.} Name the axes, write text on the chart, point at a place with an arrow, and place
the legend. In the guide: \href{https://matplotlib.org/stable/users/explain/quick_start.html}{Quick
start guide}, ``Labelling plots'';
\href{https://matplotlib.org/stable/users/explain/text/text_intro.html}{Text in Matplotlib},
\href{https://matplotlib.org/stable/users/explain/text/annotations.html}{Annotations} and the
\href{https://matplotlib.org/stable/users/explain/axes/legend_guide.html}{Legend guide}.

\textbf{In simplydata.} \code{label\_x}, \code{label\_y}, \code{write\_title\_on\_top} and
\code{write\_caption\_at\_bottom} write the text around the chart. \code{write\_note} is
\code{ax.text}, and with \code{with\_arrow=True} it is \code{ax.annotate}. The text goes to
matplotlib as you write it, so matplotlib's maths between dollar signs works too. A legend appears by
itself when \code{colour=} makes groups; \code{put\_legend\_inside}, \code{put\_legend\_at\_bottom}
and the other \code{put\_legend\_} words place it, \code{label\_legend} names it,
\code{order\_legend} orders it and \code{hide\_legend} hides it.

\begin{lstlisting}[style=py]
growth = sdata.get_example("growth")
chart = sdata.chart(growth)
chart.draw_lines(x="year", y="population", colour="country")
chart.label_y("People (millions)")
chart.label_x("Year, $t$")
chart.write_note("fastest growth", x=2020, y=12, with_arrow=True)
chart.write_caption_at_bottom("Made-up data")
chart.put_legend_inside("top left")
chart.label_legend("Country")
chart.save("labels.png")
\end{lstlisting}

\textbf{Left to matplotlib.} Arrows with your own shape and connection, text along a path, and
boxes of many styles. Hand the chart over:

\begin{lstlisting}[style=py]
growth = sdata.get_example("growth")
chart = sdata.chart(growth)
chart.draw_lines(x="year", y="population", colour="country")
fig = chart.as_matplotlib()
ax = fig.axes[0]
ax.annotate("a bend", xy=(2010, 6), xytext=(2003, 10),
            arrowprops={"arrowstyle": "-|>", "connectionstyle": "arc3,rad=0.3"},
            bbox={"boxstyle": "round", "facecolor": "lightyellow"})
fig.savefig("own_arrow.png")
\end{lstlisting}

\unsection{Scales, ticks, dates and groups on an axis}{from-matplotlib:axes}

\textbf{The job.} Choose the range and the scale of an axis (log, square root), where the ticks go and
how their numbers are written, dates and text on an axis, and a second axis. In the guide:
\href{https://matplotlib.org/stable/users/explain/quick_start.html}{Quick start guide}, ``Axis
scales and ticks'';
\href{https://matplotlib.org/stable/users/explain/axes/axes_scales.html}{Axis scales} and
\href{https://matplotlib.org/stable/users/explain/axes/axes_ticks.html}{Axis ticks}.

\textbf{In simplydata.} \code{set\_range\_x} and \code{set\_range\_y} are \code{set\_xlim} and
\code{set\_ylim}; \code{use\_log\_scale\_y} and \code{use\_scale\_y("square root")} are
\code{set\_yscale}; \code{put\_ticks\_x} is \code{set\_xticks} (with \code{labels=} it is
\code{set\_xticklabels}, and \code{every=} a \code{MultipleLocator}); \code{write\_ticks\_y\_as}
chooses a formatter; \code{rotate\_labels\_x} turns the labels; \code{flip\_y} is
\code{invert\_yaxis}. Dates and text on an axis need nothing: simplydata finds them. Months and
weekdays come in calendar order, and \code{order\_x} sets any other order. \code{add\_second\_y} is
\code{twinx}, and \code{add\_second\_y(formula=...)} is \code{secondary\_yaxis}.

\begin{lstlisting}[style=py]
shop = sdata.get_example("shop")
chart = sdata.chart(shop)
chart.draw_lines(x="date", y="sales", colour="shop")
chart.write_ticks_x_as("Jan 2026")
chart.write_ticks_y_as("thousands")
chart.rotate_labels_x(45)
chart.save("dates.png")

marks = sdata.get_example("marks")
chart = sdata.chart(marks)
chart.draw_bars(x="class", y="marks")
chart.order_x("C", "B", "A")
chart.set_range_y(0, 100)
chart.put_ticks_y(0, 25, 50, 75, 100)
chart.save("groups.png")

weather = sdata.get_example("weather")
chart = sdata.chart(weather)
chart.draw_lines(x="month", y="temperature", colour="city")
chart.add_second_y(formula="y * 9 / 5 + 32")
chart.save("second_axis.png")
\end{lstlisting}

\textbf{Left to matplotlib.} A tick formatter written as your own function, and scales of your own
(\code{FuncScale}). Hand the chart over:

\begin{lstlisting}[style=py]
from matplotlib.ticker import FuncFormatter

marks = sdata.get_example("marks")
chart = sdata.chart(marks)
chart.draw_dots(x="hours", y="marks")
fig = chart.as_matplotlib()
fig.axes[0].xaxis.set_major_formatter(FuncFormatter(lambda value, place: f"{value:.0f} h"))
fig.savefig("own_ticks.png")
\end{lstlisting}

\unsection{Numbers as colours}{from-matplotlib:colour-maps}

\textbf{The job.} Show a number as a colour (a heat map or a contour map), choose the colour map and
how numbers map to colours, and add a colour bar. In the guide:
\href{https://matplotlib.org/stable/users/explain/quick_start.html}{Quick start guide}, ``Color
mapped data''; \href{https://matplotlib.org/stable/users/explain/colors/colormaps.html}{Choosing
Colormaps} and \href{https://matplotlib.org/stable/users/explain/colors/colormapnorms.html}{Colormap
normalization}.

\textbf{In simplydata.} \code{draw\_heatmap} is \code{imshow} or \code{pcolormesh} (seaborn
\code{heatmap}), and \code{draw\_contour} is \code{contourf}, also for scattered points.
\code{use\_colour\_scale} is \code{cmap=}: it chooses the colour map. Its options are the
normalizations: \code{middle=0} puts the middle colour at 0 (a \code{TwoSlopeNorm}), \code{log=True}
spreads the colours on a log scale, \code{steps=} cuts the scale into bands, and \code{between=} sets
\code{vmin} and \code{vmax}. The colour bar comes by itself; \code{label\_colour\_bar} names it and
\code{hide\_colour\_bar} hides it.

\begin{lstlisting}[style=py]
weather = sdata.get_example("weather")
chart = sdata.chart(weather)
chart.draw_heatmap(x="month", y="city", value="rain", with_numbers=True)
chart.use_colour_scale("blues", between=[0, 300])
chart.label_colour_bar("Rain (mm)")
chart.save("heatmap.png")

island = sdata.get_example("island")
chart = sdata.chart(island)
chart.draw_contour(x="east", y="north", z="height")
chart.use_colour_scale("terrain", steps=8)
chart.save("contour.png")
\end{lstlisting}

\unsection{Several charts in one picture}{from-matplotlib:several}

\textbf{The job.} Put several Axes in one Figure: a grid of charts, one chart for each group, charts of
different sizes, or a small chart inside a big one. In the guide:
\href{https://matplotlib.org/stable/users/explain/quick_start.html}{Quick start guide}, ``Working
with multiple Figures and Axes'';
\href{https://matplotlib.org/stable/users/explain/axes/arranging_axes.html}{Arranging multiple Axes in
a Figure} and the
\href{https://matplotlib.org/stable/users/explain/axes/constrainedlayout_guide.html}{Constrained layout
guide}.

\textbf{In simplydata.} \code{draw\_one\_chart\_for\_each\_group\_in} is seaborn's
\code{FacetGrid} (\code{col=}, \code{row=}, \code{col\_wrap=}). \code{sdata.put\_side\_by\_side},
\code{sdata.put\_one\_above\_other} and \code{sdata.put\_in\_rows\_and\_columns} are
\code{plt.subplots} and \code{subplot\_mosaic} for charts you made one by one; a nested list makes
a chart span several places, and \code{with\_letters=True} writes A, B, C. \code{put\_chart\_inside} is
\code{inset\_axes}, and \code{add\_zoom\_box} is \code{indicate\_inset\_zoom}. The layout is always
matplotlib's constrained layout, so labels do not overlap.

\begin{lstlisting}[style=py]
marks = sdata.get_example("marks")
chart = sdata.chart(marks)
chart.draw_dots(x="hours", y="marks")
chart.draw_one_chart_for_each_group_in("subject", per_row=3)
chart.save("one_per_group.png")

dots = sdata.chart(marks).draw_dots(x="hours", y="marks")
bars = sdata.chart(marks).draw_bars(x="class", y="marks")
boxes = sdata.chart(marks).draw_boxes(x="subject", y="marks")
page = sdata.put_in_rows_and_columns([[dots, dots], [bars, boxes]], with_letters=True)
page.save("page.png")
\end{lstlisting}

\textbf{Left to matplotlib.} Axes placed by hand at any place in the Figure (\code{fig.add\_axes}).
\code{page.as\_matplotlib()} gives the page as one Figure: each chart is a SubFigure in
\code{fig.subfigs}, and \code{fig.axes} holds every Axes of every chart.

\unsection{Looks and settings}{from-matplotlib:style}

\textbf{The job.} Change the look of every chart at once: a style sheet, the fonts and the sizes of the
text (matplotlib's \code{rcParams}). In the guide:
\href{https://matplotlib.org/stable/users/explain/customizing.html}{Customizing Matplotlib with style
sheets and rcParams}.

\textbf{In simplydata.} \code{use\_style} picks one of five looks: \code{"paper"} (normal),
\code{"school"}, \code{"presentation"}, \code{"dark"} and \code{"black and white"}.
\code{use\_font} and \code{use\_font\_size} set the text. Each chart keeps its own look; nothing
changes matplotlib's settings for other code.

\begin{lstlisting}[style=py]
marks = sdata.get_example("marks")
chart = sdata.chart(marks)
chart.draw_boxes(x="class", y="marks")
chart.use_style("presentation")
chart.use_font("serif")
chart.use_font_size(16, title=22)
chart.save("presentation.png")
\end{lstlisting}

\textbf{Left to matplotlib.} Any other setting of \code{rcParams}. A setting such as the colour of the
axis lines is used while matplotlib draws, so change the drawn parts after the hand-over:

\begin{lstlisting}[style=py]
marks = sdata.get_example("marks")
chart = sdata.chart(marks)
chart.draw_boxes(x="class", y="marks")
fig = chart.as_matplotlib()
for side in fig.axes[0].spines.values():
    side.set_color("steelblue")
fig.savefig("own_settings.png")
\end{lstlisting}

\unsection{Lines, bands and shapes on a chart}{from-matplotlib:marks}

\textbf{The job.} Draw a reference line, shade a range, draw a rectangle, a circle or an arrow. In the
guide: \href{https://matplotlib.org/stable/plot_types/index.html}{Plot types} and
\href{https://matplotlib.org/stable/users/explain/text/annotations.html}{Annotations}.

\textbf{In simplydata.} \code{draw\_line\_at\_y} is \code{axhline} (with \code{between=} it is
\code{hlines}), \code{draw\_line\_at\_x} is \code{axvline}, \code{draw\_line\_with\_slope} is
\code{axline}, \code{shade\_between\_x} is \code{axvspan}, \code{draw\_band} is
\code{fill\_between}, \code{draw\_rectangle} and \code{draw\_circle} are the \code{Rectangle} and
\code{Circle} patches, and \code{draw\_arrow} is an arrow from one point to another.

\begin{lstlisting}[style=py]
marks = sdata.get_example("marks")
chart = sdata.chart(marks)
chart.draw_dots(x="hours", y="marks")
chart.draw_line_at_y(50, labelled="pass mark")
chart.shade_between_x(6, 8)
chart.draw_rectangle(x_between=[3, 5], y_between=[60, 80])
chart.draw_arrow(start=[2, 40], end=[4, 60])
chart.save("shapes.png")
\end{lstlisting}

\textbf{Left to matplotlib.} Polygons, paths and other patches. Hand the chart over and add them:

\begin{lstlisting}[style=py]
from matplotlib.patches import Polygon

marks = sdata.get_example("marks")
chart = sdata.chart(marks)
chart.draw_dots(x="hours", y="marks")
fig = chart.as_matplotlib()
fig.axes[0].add_patch(Polygon([[2, 40], [5, 90], [8, 60]], alpha=0.2, color="purple"))
fig.savefig("polygon.png")
\end{lstlisting}

\unsection{Charts in 3D}{from-matplotlib:3d}

\textbf{The job.} Draw a surface, dots or lines in three dimensions and choose the view. In the guide:
\href{https://matplotlib.org/stable/users/explain/toolkits/mplot3d.html}{The mplot3d toolkit}.

\textbf{In simplydata.} \code{draw\_surface} is \code{plot\_surface} (with \code{as\_wires=True} it is
\code{plot\_wireframe}), \code{draw\_dots\_3d} is \code{scatter} on 3D Axes, \code{draw\_lines\_3d}
is \code{plot} in 3D and \code{draw\_bars\_3d} is \code{bar3d}. \code{look\_from} is
\code{view\_init}, and \code{label\_z}, \code{set\_range\_z} and \code{put\_ticks\_z} work on the third
axis.

\begin{lstlisting}[style=py]
island = sdata.get_example("island")
chart = sdata.chart(island)
chart.draw_surface(x="east", y="north", z="height")
chart.look_from(height=40, around=30)
chart.label_z("Height (m)")
chart.save("surface.png")
\end{lstlisting}

\unsection{Moving charts}{from-matplotlib:animations}

\textbf{The job.} Make a chart that moves, one picture after another. In the guide:
\href{https://matplotlib.org/stable/users/explain/animations/animations.html}{Animations using
Matplotlib}.

\textbf{In simplydata.} \code{animate(over="year")} draws one picture for each value of a column:
matplotlib draws each picture and Pillow writes them into a .gif file.

\begin{lstlisting}[style=py]
growth = sdata.get_example("growth")
chart = sdata.chart(growth)
chart.draw_bars(x="country", y="population")
chart.animate(over="year")
chart.save("growth.gif")
\end{lstlisting}

\unsection{Saving, showing and web pages}{from-matplotlib:saving}

\textbf{The job.} Save the picture in a file type and size you need, show it in a window, or make a
chart you can point at in a web page. In the guide:
\href{https://matplotlib.org/stable/users/explain/figure/index.html}{Figures and backends}, with
\href{https://matplotlib.org/stable/users/explain/figure/backends.html}{Backends} and
\href{https://matplotlib.org/stable/users/explain/figure/interactive.html}{Interactive figures}.

\textbf{In simplydata.} \code{save} is \code{savefig}: the end of the file name picks the type
(.png, .pdf, .svg, .eps, .tif, .jpg, .gif or .html), the picture is 300 dots per inch and trimmed to
what is drawn, \code{resolution=} is \code{dpi}, and \code{see\_through=True} is
\code{transparent=True}. \code{set\_size} is \code{figsize}, in inches or cm. \code{show} is
\code{plt.show}. A .html file and \code{show\_in\_web\_page} make a web page, drawn by plotly from
what matplotlib drew; \code{show\_when\_pointing} chooses what it shows when the pointer rests on a
dot. In Jupyter a chart shows by itself.

\begin{lstlisting}[style=py]
marks = sdata.get_example("marks")
chart = sdata.chart(marks)
chart.draw_dots(x="hours", y="marks", colour="class")
chart.set_size(width=17, height=10, unit="cm")
chart.save("for_the_journal.pdf")
chart.save("for_a_slide.png", resolution=150, see_through=True)
chart.get_size(unit="cm")
\end{lstlisting}

% needs: plotly
\begin{lstlisting}[style=py]
marks = sdata.get_example("marks")
chart = sdata.chart(marks)
chart.draw_dots(x="hours", y="marks", colour="class")
chart.show_when_pointing("name", "subject")
chart.save("marks.html")
\end{lstlisting}

\textbf{The hand-over and web pages.} \code{chart.as\_matplotlib()} always gives the matplotlib
Figure, also for a chart you save as a web page: the web page is a copy that plotly draws from it.
Changes you make in matplotlib after the hand-over are not in \code{chart.save("marks.html")}.

\textbf{Left to matplotlib.} Choosing the backend, windows that react to the mouse or the keyboard
(events), and widgets such as sliders and buttons.

\unsection{Relationships: dots and lines}{from-matplotlib:relational}

\textbf{The job.} Show how one number goes with another: dots for each row, or a line through the
average at each x, with groups by colour, size or shape. In the tutorial:
\href{https://seaborn.pydata.org/tutorial/relational.html}{Visualizing statistical relationships}
(\code{relplot}, \code{scatterplot}, \code{lineplot}), and the overview of the
\href{https://seaborn.pydata.org/tutorial/function_overview.html}{plotting functions} and the
\href{https://seaborn.pydata.org/tutorial/data_structure.html}{data structures} seaborn takes.

\textbf{In simplydata.} \code{draw\_dots} is \code{scatterplot}: \code{colour=} is \code{hue=},
\code{size=} is \code{size=} and \code{shape=} is \code{style=}. \code{draw\_lines} is
\code{lineplot}; it draws the average of the rows at each x, and \code{error=} adds the band.
\code{relplot}'s panels are \code{draw\_one\_chart\_for\_each\_group\_in}.

\begin{lstlisting}[style=py]
marks = sdata.get_example("marks")
chart = sdata.chart(marks)
chart.draw_dots(x="hours", y="marks", colour="class", shape="subject")
chart.save("relationship.png")

school = sdata.get_example("school")
chart = sdata.chart(school)
chart.draw_lines(x="sleep", y="marks", colour="breakfast", error="confidence interval")
chart.save("average_line.png")
\end{lstlisting}

\unsection{Distributions}{from-matplotlib:distributions}

\textbf{The job.} Show how the values of a column are spread: a histogram, a smooth density, the
share of values up to each point, a rug, and two columns together. In the tutorial:
\href{https://seaborn.pydata.org/tutorial/distributions.html}{Visualizing distributions of data}
(\code{displot}, \code{histplot}, \code{kdeplot}, \code{ecdfplot}, \code{rugplot}).

\textbf{In simplydata.} \code{draw\_histogram} is \code{histplot} (\code{with\_bins=},
\code{show="density"}, \code{stacked=True}, \code{cumulative=True}); \code{draw\_density} is
\code{kdeplot}; \code{draw\_cumulative\_share} is \code{ecdfplot}; \code{draw\_rug\_lines} is
\code{rugplot}. With \code{y=} as well, \code{draw\_histogram} and \code{draw\_density} show two
columns together, and \code{draw\_side\_histograms} adds the edges of \code{jointplot}.
\code{fit\_distribution} draws a fitted curve over a histogram.

\begin{lstlisting}[style=py]
marks = sdata.get_example("marks")
chart = sdata.chart(marks)
chart.draw_histogram(x="marks", colour="class", show="density")
chart.draw_density(x="marks", colour="class")
chart.draw_rug_lines(x="marks")
chart.save("spread.png")

chart = sdata.chart(marks)
chart.draw_dots(x="hours", y="marks")
chart.draw_side_histograms(drawn_as="density")
chart.save("joint.png")
\end{lstlisting}

\unsection{Groups: categorical charts}{from-matplotlib:categorical}

\textbf{The job.} Compare the values of groups: every dot, boxes, violins, bars of an average with
an interval, points joined across groups, and counts. In the tutorial:
\href{https://seaborn.pydata.org/tutorial/categorical.html}{Visualizing categorical data}
(\code{catplot}, \code{stripplot}, \code{swarmplot}, \code{boxplot}, \code{violinplot},
\code{barplot}, \code{pointplot}, \code{countplot}).

\textbf{In simplydata.} \code{draw\_all\_dots} is \code{stripplot} and \code{swarmplot};
\code{draw\_boxes} is \code{boxplot}; \code{draw\_violins} is \code{violinplot} (\code{in\_halves=}
is \code{split=True}); \code{draw\_bars} is \code{barplot}, with \code{show=} for the average,
median, total or count; \code{draw\_interaction\_lines} is \code{pointplot}; \code{draw\_count\_bars}
is \code{countplot}. \code{catplot}'s panels are \code{draw\_one\_chart\_for\_each\_group\_in}.

\begin{lstlisting}[style=py]
marks = sdata.get_example("marks")
chart = sdata.chart(marks)
chart.draw_boxes(x="subject", y="marks")
chart.draw_all_dots(x="subject", y="marks")
chart.save("boxes_and_dots.png")

school = sdata.get_example("school")
chart = sdata.chart(school)
chart.draw_interaction_lines(x="class", y="marks", colour="breakfast")
chart.save("points.png")

chart = sdata.chart(school)
chart.draw_count_bars(x="class", colour="breakfast")
chart.save("counts.png")
\end{lstlisting}

\textbf{Left to seaborn.} \code{boxenplot}, the letter-value boxes for big tables. Hand an empty chart
over and let seaborn draw on its Axes:

\begin{lstlisting}[style=py]
import seaborn as sns

marks = sdata.get_example("marks")
chart = sdata.chart(marks)
chart.label_x("subject")
fig = chart.as_matplotlib()
sns.boxenplot(marks.as_pandas(), x="subject", y="marks", ax=fig.axes[0])
fig.savefig("boxen.png")
\end{lstlisting}

\unsection{Averages with error bars}{from-matplotlib:error-bars}

\textbf{The job.} Show how sure an average is (a confidence interval or a standard error) or how
spread the values are (a standard deviation). In the tutorial:
\href{https://seaborn.pydata.org/tutorial/error_bars.html}{Statistical estimation and error bars}.

\textbf{In simplydata.} \code{error=} on \code{draw\_bars}, \code{draw\_lines} and \code{draw\_dots}
takes \code{"confidence interval"}, \code{"standard error"} or \code{"standard deviation"}, as
seaborn's \code{errorbar=("ci", 95)}, \code{"se"} and \code{"sd"}. \code{draw\_error\_bars} adds them
to dots or lines. The interval is the t interval, not seaborn's bootstrap, so the same table always
gives the same bars, and a line under the chart says which one was drawn.

\begin{lstlisting}[style=py]
marks = sdata.get_example("marks")
chart = sdata.chart(marks)
chart.draw_bars(x="class", y="marks", show="average", error="confidence interval")
chart.save("error_bars.png")
\end{lstlisting}

\unsection{Lines that fit}{from-matplotlib:regression}

\textbf{The job.} Draw the line that fits the dots best, with its confidence band, a curve, a smooth
line, and what is left over. In the tutorial:
\href{https://seaborn.pydata.org/tutorial/regression.html}{Estimating regression fits}
(\code{regplot}, \code{lmplot}, \code{residplot}).

\textbf{In simplydata.} \code{fit\_line} is \code{regplot}; \code{show\_equation=True} writes the
equation and $R^2$ on the chart. \code{fit\_curve(of\_degree=2)} is \code{order=2}.
\code{draw\_smooth\_line} is \code{lowess=True}. \code{draw\_residual\_dots} is \code{residplot}.
For \code{lmplot}'s panels, add \code{draw\_one\_chart\_for\_each\_group\_in}.

\begin{lstlisting}[style=py]
marks = sdata.get_example("marks")
chart = sdata.chart(marks)
chart.draw_dots(x="hours", y="marks")
chart.fit_line(show_equation=True)
chart.save("fitted_line.png")

chart = sdata.chart(marks)
chart.draw_residual_dots(x="hours", y="marks")
chart.save("left_over.png")
\end{lstlisting}

\unsection{Grids of charts}{from-matplotlib:grids}

\textbf{The job.} One chart for each group in a grid, every column against every other, and a chart
with its edges. In the tutorial: \href{https://seaborn.pydata.org/tutorial/axis_grids.html}{Building
structured multi-plot grids} (\code{FacetGrid}, \code{PairGrid}, \code{pairplot}, \code{JointGrid},
\code{jointplot}).

\textbf{In simplydata.} \code{draw\_one\_chart\_for\_each\_group\_in} with \code{down=} and
\code{across=} is \code{FacetGrid} with \code{row=} and \code{col=}; \code{same\_y=False} is
\code{sharey=False}. \code{sdata.draw\_all\_pairs} is \code{pairplot}. \code{draw\_side\_histograms}
is \code{jointplot}.

\begin{lstlisting}[style=py]
marks = sdata.get_example("marks")
chart = sdata.chart(marks)
chart.draw_dots(x="hours", y="marks")
chart.draw_one_chart_for_each_group_in(down="class", across="subject")
chart.save("grid.png")

pairs = sdata.draw_all_pairs(marks, colour="class")
pairs.save("pairs.png")
\end{lstlisting}

\textbf{Left to seaborn.} Your own function mapped on each panel of a grid (\code{FacetGrid.map}), and
different kinds of charts above and below the diagonal of a \code{PairGrid}.

\unsection{The look of a chart}{from-matplotlib:aesthetics}

\textbf{The job.} Choose the overall look, take away lines of the frame, and scale everything for a
paper, a talk or a poster. In the tutorial:
\href{https://seaborn.pydata.org/tutorial/aesthetics.html}{Controlling figure aesthetics}
(\code{set\_theme}, \code{set\_style}, \code{despine}, \code{set\_context}).

\textbf{In simplydata.} The normal look (\code{"paper"}) already leaves out the top and right lines,
as \code{despine} does. \code{use\_style} is \code{set\_style} and \code{set\_context} in one word:
\code{"presentation"} is like seaborn's \code{"talk"}. \code{hide\_frame} takes the whole frame away,
\code{show\_grid} and \code{hide\_grid} the grid lines, and \code{use\_background\_colour} the
background.

\begin{lstlisting}[style=py]
weather = sdata.get_example("weather")
chart = sdata.chart(weather)
chart.draw_lines(x="month", y="rain", colour="city")
chart.use_style("dark")
chart.show_grid()
chart.save("dark.png")
\end{lstlisting}

\unsection{Colour palettes}{from-matplotlib:palettes}

\textbf{The job.} Choose colours for groups (a qualitative palette), for numbers that go from low to
high (sequential) and for numbers on both sides of a middle (diverging). In the tutorial:
\href{https://seaborn.pydata.org/tutorial/color_palettes.html}{Choosing color palettes}.

\textbf{In simplydata.} The normal colours for groups are the Okabe-Ito set, safe for readers who do
not see some colours. \code{use\_colours} is \code{color\_palette} and \code{set\_palette}: a named set
(\code{"Set2"}, \code{"deep"}, \code{"colorblind"}), your own colours, or \code{how\_many=} colours
from a colour map. \code{use\_colour\_scale} is the \code{cmap} of a heat map; \code{"blue to red"}
with \code{middle=0} is a diverging palette centred on 0.

\begin{lstlisting}[style=py]
marks = sdata.get_example("marks")
chart = sdata.chart(marks)
chart.draw_bars(x="class", y="marks", colour="subject")
chart.use_colours("viridis", how_many=3)
chart.save("palette.png")
\end{lstlisting}

\unsection{The jobs left to matplotlib and seaborn}{from-matplotlib:left-out}

simplydata draws the charts most people need in research and teaching, with one word each. These jobs
are left out on purpose. Each one is a hand-over: \code{fig = chart.as\_matplotlib()}, then
matplotlib or seaborn code on \code{fig.axes}, then \code{fig.savefig(...)}.
\begin{itemize}
  \item \textbf{Arrows and vector fields} (\code{quiver}, \code{barbs}, \code{streamplot}).
  \item \textbf{Polygons, paths and other patches} (\code{fill}, \code{Polygon}), shown above.
  \item \textbf{A table drawn inside a chart} (\code{ax.table}).
  \item \textbf{Text of your own shape}: arrow styles, text along a path, many box styles (shown above).
  \item \textbf{Tick formatters as your own function} and scales of your own (shown above).
  \item \textbf{Axes placed by hand} (\code{fig.add\_axes}) and margins (\code{ax.margins}).
  \item \textbf{Images} (\code{imread}, \code{imsave}, \code{figimage}); images get their own words in
    a later phase.
  \item \textbf{Interaction}: backends, mouse and key events, widgets (\code{Slider}, \code{Button}).
  \item \textbf{More 3D}: voxels, arrows and contour sets in 3D.
  \item \textbf{seaborn's letter-value boxes} (\code{boxenplot}, shown above), a function of your own on
    each panel of a grid (\code{FacetGrid.map}) and the objects interface (\code{seaborn.objects}).
  \item \textbf{Any \code{rcParams} setting} beyond the looks, fonts and sizes of \code{use\_style},
    \code{use\_font} and \code{use\_font\_size} (shown above).
\end{itemize}

One more example: a vector field over a contour map.

\begin{lstlisting}[style=py]
island = sdata.get_example("island")
chart = sdata.chart(island)
chart.draw_contour(x="east", y="north", z="height")
fig = chart.as_matplotlib()
ax = fig.axes[0]
ax.quiver([2000, 4000], [1000, 3000], [1, -1], [1, 0.5], color="white")
fig.savefig("vectors.png")
\end{lstlisting}
).
% The topics follow the order of the matplotlib Quick start guide and User guide, then the seaborn tutorial, used as a
% checklist only: the text and the examples are our own, and each section links to the official page. Every example in
% a py box is run by tests/test_from_matplotlib_chapter.py (with "import simplydata as sdata" in front, and matplotlib's
% Agg backend), so it cannot go out of date. A line "% needs: plotly" just before a box says the box needs that package.

\unchapter{From matplotlib and seaborn to simplydata}{from-matplotlib}

This chapter is for you if you already know matplotlib or seaborn, or if you want to learn them after
simplydata. It walks through the topics of the official guides, in their order: the
\href{https://matplotlib.org/stable/users/explain/quick_start.html}{matplotlib Quick start guide} and
\href{https://matplotlib.org/stable/users/index.html}{User guide}, then the
\href{https://seaborn.pydata.org/tutorial.html}{seaborn tutorial}. For each topic it says what the job
is, which simplydata words do it, and a short example on the practice tables. Some jobs are left to
matplotlib and seaborn on purpose. For those, the chapter shows how to hand the chart over and keep
going in matplotlib.

Every example starts with \code{import simplydata as sdata}. It is written once here, not in each
box. The examples use the practice tables of \code{sdata.get\_example}, so you can run them as they
are.

\unsection{Five things to know first}{from-matplotlib:first}

\begin{enumerate}
  \item \textbf{A chart is a list of steps.} In matplotlib each call draws at once. In simplydata,
    \code{chart.draw\_dots(...)} and \code{chart.label\_x(...)} only write down a step. The chart is
    drawn when you save it, show it, or look at it in Jupyter. So the order of the steps rarely
    matters, and \code{chart.undo()} takes the last step back.
  \item \textbf{There is no figure and no axes to hold.} simplydata makes the matplotlib Figure and
    its Axes for you each time it draws. You never write \code{fig, ax = plt.subplots()}.
  \item \textbf{Columns by name.} You give the table once, to \code{sdata.chart(table)}, and then the
    column names, as seaborn does with \code{data=}. A word works on a pandas or a Polars table.
  \item \textbf{One word for one job.} matplotlib often has several ways to do the same thing
    (\code{plt.title}, \code{ax.set\_title}, \code{fig.suptitle}). simplydata has one word,
    \code{write\_title\_on\_top}, and \code{sdata.help("ax.set\_title")} names it.
  \item \textbf{The door is always open.} \code{fig = chart.as\_matplotlib()} gives the chart as a
    plain matplotlib Figure. Its charts are the Axes in \code{fig.axes}, so \code{ax = fig.axes[0]}.
    Then any matplotlib code works. It is a copy: changing \code{fig} does not change the chart.
    Save it with \code{fig.savefig(...)}.
\end{enumerate}

\begin{lstlisting}[style=py]
marks = sdata.get_example("marks")
chart = sdata.chart(marks)
chart.draw_dots(x="hours", y="marks")
fig = chart.as_matplotlib()              # a matplotlib Figure: use any matplotlib code here
ax = fig.axes[0]
ax.set_facecolor("whitesmoke")
fig.savefig("handed_over.png", dpi=150)
\end{lstlisting}

Do you know a matplotlib or seaborn name, but not the simplydata word? Ask:
\code{sdata.help("ax.set\_xlabel")}, \code{sdata.help("plt.hist")} or
\code{sdata.help("sns.lineplot")}. The name can have \code{plt.}, \code{ax.}, \code{fig.} or
\code{sns.} in front. The full list is in the section ``Coming from another tool?'' of the simplydata
dictionary
(\href{https://github.com/bioforkgithub/simplydata/blob/main/docs/manual/simplydataDictionary.pdf}%
{simplydataDictionary.pdf}).

\unsection{A first chart}{from-matplotlib:first-chart}

\textbf{The job.} Draw a line from two columns, give it a title and axis names, and save it as a
picture. In the guide: \href{https://matplotlib.org/stable/users/explain/quick_start.html}{Quick
start guide}, the simple example at its start.

\textbf{In simplydata.} \code{draw\_lines} is \code{ax.plot}; \code{write\_title\_on\_top},
\code{label\_x} and \code{label\_y} are \code{set\_title}, \code{set\_xlabel} and
\code{set\_ylabel}; \code{save} is \code{fig.savefig}. Without \code{label\_x}, the column name is the
axis name.

\begin{lstlisting}[style=py]
weather = sdata.get_example("weather")
chart = sdata.chart(weather)
chart.draw_lines(x="month", y="temperature", colour="city")
chart.write_title_on_top("Temperature through the year")
chart.label_y("Temperature (degrees C)")
chart.save("first_chart.png")
\end{lstlisting}

\unsection{The parts of a figure}{from-matplotlib:parts}

\textbf{The job.} Know the parts matplotlib draws: the Figure (the whole picture), the Axes (one
chart inside it), each Axis (the x and y scales) and the Artists (every line, dot, bar and piece of
text). In the guide: \href{https://matplotlib.org/stable/users/explain/quick_start.html}{Quick start
guide}, ``Parts of a Figure'', and
\href{https://matplotlib.org/stable/users/explain/artists/index.html}{Artists}.

\textbf{In simplydata.} A simplydata chart is one Figure with one Axes, or several Axes after
\code{draw\_one\_chart\_for\_each\_group\_in}. You do not touch the parts. You can read them back as
tables: \code{get\_numbers\_drawn} reads the lines, dots and bars, \code{get\_text\_drawn} the text,
and \code{get\_scales\_drawn} each axis. \code{as\_matplotlib} gives the parts themselves.

\begin{lstlisting}[style=py]
marks = sdata.get_example("marks")
chart = sdata.chart(marks)
chart.draw_bars(x="class", y="marks")
chart.get_numbers_drawn()                # one row for each bar: where it is and how high
chart.get_scales_drawn()                 # the range and kind of each axis
fig = chart.as_matplotlib()
print(len(fig.axes), "Axes in the Figure")
\end{lstlisting}

\unsection{The data a chart takes}{from-matplotlib:data}

\textbf{The job.} Give the numbers to draw. matplotlib takes lists or NumPy arrays, or column names
with \code{data=}. In the guide: \href{https://matplotlib.org/stable/users/explain/quick_start.html}%
{Quick start guide}, ``Types of inputs to plotting functions''.

\textbf{In simplydata.} A table and column names is the usual way: \code{sdata.chart(table)} reads a
table, a file name or a pandas or Polars DataFrame. Without a table, give lists to the step itself.
Numbers made with \code{sdata.number\_list} start their own chart.

\begin{lstlisting}[style=py]
chart = sdata.chart()                    # no table: lists go to the step
chart.draw_lines(x=[1, 2, 3, 4], y=[10, 12, 9, 14])
chart.save("from_lists.png")

numbers = sdata.number_list([4, 8, 15, 16, 23, 42])
chart = numbers.start_chart()
chart.draw_histogram(x="value")
chart.save("from_numbers.png")
\end{lstlisting}

\unsection{Two ways of writing, one in simplydata}{from-matplotlib:styles}

\textbf{The job.} matplotlib can be written in two styles: the explicit one (\code{fig, ax =
plt.subplots()}, then \code{ax.plot(...)}) and the implicit one with \code{pyplot}
(\code{plt.plot(...)}, which draws on the current Axes). In the guide:
\href{https://matplotlib.org/stable/users/explain/quick_start.html}{Quick start guide}, ``Coding
styles''.

\textbf{In simplydata.} There is one style. The chart is the object you name, as in the explicit style,
and nothing is ``current''. So two charts never draw on each other. To make many charts the same way,
write a small function that gives back a chart, as the guide suggests for matplotlib.

\begin{lstlisting}[style=py]
def marks_chart(subject):
    marks = sdata.get_example("marks")
    marks.keep_rows_where("subject", is_equal_to=subject)
    chart = sdata.chart(marks)
    chart.draw_dots(x="hours", y="marks")
    chart.write_title_on_top(subject)
    return chart


marks_chart("Maths").save("maths.png")
marks_chart("Science").save("science.png")
\end{lstlisting}

\unsection{Colours, lines and dots}{from-matplotlib:styling}

\textbf{The job.} Choose the colour, width and dash of a line, and the size, shape and transparency of
dots. In the guide: \href{https://matplotlib.org/stable/users/explain/quick_start.html}{Quick start
guide}, ``Styling Artists'', and
\href{https://matplotlib.org/stable/users/explain/colors/colors.html}{Specifying colors}.

\textbf{In simplydata.} Options with plain names on the draw words: \code{thickness=} is
\code{linewidth}, \code{dashed=True} and \code{dotted=True} are \code{linestyle}, \code{size=} is the
dot size (a column, for sizes that change), \code{transparency=} is one minus \code{alpha}.
\code{use\_colour} sets one colour, \code{use\_colours} one colour for each group, and
\code{use\_shapes} a shape for each group. The groups come from \code{colour="column"}, as seaborn's
\code{hue=}.

\begin{lstlisting}[style=py]
weather = sdata.get_example("weather")
chart = sdata.chart(weather)
chart.draw_lines(x="month", y="temperature", colour="city", thickness=3, dashed="city")
chart.use_colours("black", "darkorange")
chart.save("styled_lines.png")

marks = sdata.get_example("marks")
chart = sdata.chart(marks)
chart.draw_dots(x="hours", y="marks", colour="class", size="height", transparency=0.3)
chart.use_shapes("circle", "square", "triangle")
chart.save("styled_dots.png")
\end{lstlisting}

\unsection{Titles, labels, notes and legends}{from-matplotlib:labels}

\textbf{The job.} Name the axes, write text on the chart, point at a place with an arrow, and place
the legend. In the guide: \href{https://matplotlib.org/stable/users/explain/quick_start.html}{Quick
start guide}, ``Labelling plots'';
\href{https://matplotlib.org/stable/users/explain/text/text_intro.html}{Text in Matplotlib},
\href{https://matplotlib.org/stable/users/explain/text/annotations.html}{Annotations} and the
\href{https://matplotlib.org/stable/users/explain/axes/legend_guide.html}{Legend guide}.

\textbf{In simplydata.} \code{label\_x}, \code{label\_y}, \code{write\_title\_on\_top} and
\code{write\_caption\_at\_bottom} write the text around the chart. \code{write\_note} is
\code{ax.text}, and with \code{with\_arrow=True} it is \code{ax.annotate}. The text goes to
matplotlib as you write it, so matplotlib's maths between dollar signs works too. A legend appears by
itself when \code{colour=} makes groups; \code{put\_legend\_inside}, \code{put\_legend\_at\_bottom}
and the other \code{put\_legend\_} words place it, \code{label\_legend} names it,
\code{order\_legend} orders it and \code{hide\_legend} hides it.

\begin{lstlisting}[style=py]
growth = sdata.get_example("growth")
chart = sdata.chart(growth)
chart.draw_lines(x="year", y="population", colour="country")
chart.label_y("People (millions)")
chart.label_x("Year, $t$")
chart.write_note("fastest growth", x=2020, y=12, with_arrow=True)
chart.write_caption_at_bottom("Made-up data")
chart.put_legend_inside("top left")
chart.label_legend("Country")
chart.save("labels.png")
\end{lstlisting}

\textbf{Left to matplotlib.} Arrows with your own shape and connection, text along a path, and
boxes of many styles. Hand the chart over:

\begin{lstlisting}[style=py]
growth = sdata.get_example("growth")
chart = sdata.chart(growth)
chart.draw_lines(x="year", y="population", colour="country")
fig = chart.as_matplotlib()
ax = fig.axes[0]
ax.annotate("a bend", xy=(2010, 6), xytext=(2003, 10),
            arrowprops={"arrowstyle": "-|>", "connectionstyle": "arc3,rad=0.3"},
            bbox={"boxstyle": "round", "facecolor": "lightyellow"})
fig.savefig("own_arrow.png")
\end{lstlisting}

\unsection{Scales, ticks, dates and groups on an axis}{from-matplotlib:axes}

\textbf{The job.} Choose the range and the scale of an axis (log, square root), where the ticks go and
how their numbers are written, dates and text on an axis, and a second axis. In the guide:
\href{https://matplotlib.org/stable/users/explain/quick_start.html}{Quick start guide}, ``Axis
scales and ticks'';
\href{https://matplotlib.org/stable/users/explain/axes/axes_scales.html}{Axis scales} and
\href{https://matplotlib.org/stable/users/explain/axes/axes_ticks.html}{Axis ticks}.

\textbf{In simplydata.} \code{set\_range\_x} and \code{set\_range\_y} are \code{set\_xlim} and
\code{set\_ylim}; \code{use\_log\_scale\_y} and \code{use\_scale\_y("square root")} are
\code{set\_yscale}; \code{put\_ticks\_x} is \code{set\_xticks} (with \code{labels=} it is
\code{set\_xticklabels}, and \code{every=} a \code{MultipleLocator}); \code{write\_ticks\_y\_as}
chooses a formatter; \code{rotate\_labels\_x} turns the labels; \code{flip\_y} is
\code{invert\_yaxis}. Dates and text on an axis need nothing: simplydata finds them. Months and
weekdays come in calendar order, and \code{order\_x} sets any other order. \code{add\_second\_y} is
\code{twinx}, and \code{add\_second\_y(formula=...)} is \code{secondary\_yaxis}.

\begin{lstlisting}[style=py]
shop = sdata.get_example("shop")
chart = sdata.chart(shop)
chart.draw_lines(x="date", y="sales", colour="shop")
chart.write_ticks_x_as("Jan 2026")
chart.write_ticks_y_as("thousands")
chart.rotate_labels_x(45)
chart.save("dates.png")

marks = sdata.get_example("marks")
chart = sdata.chart(marks)
chart.draw_bars(x="class", y="marks")
chart.order_x("C", "B", "A")
chart.set_range_y(0, 100)
chart.put_ticks_y(0, 25, 50, 75, 100)
chart.save("groups.png")

weather = sdata.get_example("weather")
chart = sdata.chart(weather)
chart.draw_lines(x="month", y="temperature", colour="city")
chart.add_second_y(formula="y * 9 / 5 + 32")
chart.save("second_axis.png")
\end{lstlisting}

\textbf{Left to matplotlib.} A tick formatter written as your own function, and scales of your own
(\code{FuncScale}). Hand the chart over:

\begin{lstlisting}[style=py]
from matplotlib.ticker import FuncFormatter

marks = sdata.get_example("marks")
chart = sdata.chart(marks)
chart.draw_dots(x="hours", y="marks")
fig = chart.as_matplotlib()
fig.axes[0].xaxis.set_major_formatter(FuncFormatter(lambda value, place: f"{value:.0f} h"))
fig.savefig("own_ticks.png")
\end{lstlisting}

\unsection{Numbers as colours}{from-matplotlib:colour-maps}

\textbf{The job.} Show a number as a colour (a heat map or a contour map), choose the colour map and
how numbers map to colours, and add a colour bar. In the guide:
\href{https://matplotlib.org/stable/users/explain/quick_start.html}{Quick start guide}, ``Color
mapped data''; \href{https://matplotlib.org/stable/users/explain/colors/colormaps.html}{Choosing
Colormaps} and \href{https://matplotlib.org/stable/users/explain/colors/colormapnorms.html}{Colormap
normalization}.

\textbf{In simplydata.} \code{draw\_heatmap} is \code{imshow} or \code{pcolormesh} (seaborn
\code{heatmap}), and \code{draw\_contour} is \code{contourf}, also for scattered points.
\code{use\_colour\_scale} is \code{cmap=}: it chooses the colour map. Its options are the
normalizations: \code{middle=0} puts the middle colour at 0 (a \code{TwoSlopeNorm}), \code{log=True}
spreads the colours on a log scale, \code{steps=} cuts the scale into bands, and \code{between=} sets
\code{vmin} and \code{vmax}. The colour bar comes by itself; \code{label\_colour\_bar} names it and
\code{hide\_colour\_bar} hides it.

\begin{lstlisting}[style=py]
weather = sdata.get_example("weather")
chart = sdata.chart(weather)
chart.draw_heatmap(x="month", y="city", value="rain", with_numbers=True)
chart.use_colour_scale("blues", between=[0, 300])
chart.label_colour_bar("Rain (mm)")
chart.save("heatmap.png")

island = sdata.get_example("island")
chart = sdata.chart(island)
chart.draw_contour(x="east", y="north", z="height")
chart.use_colour_scale("terrain", steps=8)
chart.save("contour.png")
\end{lstlisting}

\unsection{Several charts in one picture}{from-matplotlib:several}

\textbf{The job.} Put several Axes in one Figure: a grid of charts, one chart for each group, charts of
different sizes, or a small chart inside a big one. In the guide:
\href{https://matplotlib.org/stable/users/explain/quick_start.html}{Quick start guide}, ``Working
with multiple Figures and Axes'';
\href{https://matplotlib.org/stable/users/explain/axes/arranging_axes.html}{Arranging multiple Axes in
a Figure} and the
\href{https://matplotlib.org/stable/users/explain/axes/constrainedlayout_guide.html}{Constrained layout
guide}.

\textbf{In simplydata.} \code{draw\_one\_chart\_for\_each\_group\_in} is seaborn's
\code{FacetGrid} (\code{col=}, \code{row=}, \code{col\_wrap=}). \code{sdata.put\_side\_by\_side},
\code{sdata.put\_one\_above\_other} and \code{sdata.put\_in\_rows\_and\_columns} are
\code{plt.subplots} and \code{subplot\_mosaic} for charts you made one by one; a nested list makes
a chart span several places, and \code{with\_letters=True} writes A, B, C. \code{put\_chart\_inside} is
\code{inset\_axes}, and \code{add\_zoom\_box} is \code{indicate\_inset\_zoom}. The layout is always
matplotlib's constrained layout, so labels do not overlap.

\begin{lstlisting}[style=py]
marks = sdata.get_example("marks")
chart = sdata.chart(marks)
chart.draw_dots(x="hours", y="marks")
chart.draw_one_chart_for_each_group_in("subject", per_row=3)
chart.save("one_per_group.png")

dots = sdata.chart(marks).draw_dots(x="hours", y="marks")
bars = sdata.chart(marks).draw_bars(x="class", y="marks")
boxes = sdata.chart(marks).draw_boxes(x="subject", y="marks")
page = sdata.put_in_rows_and_columns([[dots, dots], [bars, boxes]], with_letters=True)
page.save("page.png")
\end{lstlisting}

\textbf{Left to matplotlib.} Axes placed by hand at any place in the Figure (\code{fig.add\_axes}).
\code{page.as\_matplotlib()} gives the page as one Figure: each chart is a SubFigure in
\code{fig.subfigs}, and \code{fig.axes} holds every Axes of every chart.

\unsection{Looks and settings}{from-matplotlib:style}

\textbf{The job.} Change the look of every chart at once: a style sheet, the fonts and the sizes of the
text (matplotlib's \code{rcParams}). In the guide:
\href{https://matplotlib.org/stable/users/explain/customizing.html}{Customizing Matplotlib with style
sheets and rcParams}.

\textbf{In simplydata.} \code{use\_style} picks one of five looks: \code{"paper"} (normal),
\code{"school"}, \code{"presentation"}, \code{"dark"} and \code{"black and white"}.
\code{use\_font} and \code{use\_font\_size} set the text. Each chart keeps its own look; nothing
changes matplotlib's settings for other code.

\begin{lstlisting}[style=py]
marks = sdata.get_example("marks")
chart = sdata.chart(marks)
chart.draw_boxes(x="class", y="marks")
chart.use_style("presentation")
chart.use_font("serif")
chart.use_font_size(16, title=22)
chart.save("presentation.png")
\end{lstlisting}

\textbf{Left to matplotlib.} Any other setting of \code{rcParams}. A setting such as the colour of the
axis lines is used while matplotlib draws, so change the drawn parts after the hand-over:

\begin{lstlisting}[style=py]
marks = sdata.get_example("marks")
chart = sdata.chart(marks)
chart.draw_boxes(x="class", y="marks")
fig = chart.as_matplotlib()
for side in fig.axes[0].spines.values():
    side.set_color("steelblue")
fig.savefig("own_settings.png")
\end{lstlisting}

\unsection{Lines, bands and shapes on a chart}{from-matplotlib:marks}

\textbf{The job.} Draw a reference line, shade a range, draw a rectangle, a circle or an arrow. In the
guide: \href{https://matplotlib.org/stable/plot_types/index.html}{Plot types} and
\href{https://matplotlib.org/stable/users/explain/text/annotations.html}{Annotations}.

\textbf{In simplydata.} \code{draw\_line\_at\_y} is \code{axhline} (with \code{between=} it is
\code{hlines}), \code{draw\_line\_at\_x} is \code{axvline}, \code{draw\_line\_with\_slope} is
\code{axline}, \code{shade\_between\_x} is \code{axvspan}, \code{draw\_band} is
\code{fill\_between}, \code{draw\_rectangle} and \code{draw\_circle} are the \code{Rectangle} and
\code{Circle} patches, and \code{draw\_arrow} is an arrow from one point to another.

\begin{lstlisting}[style=py]
marks = sdata.get_example("marks")
chart = sdata.chart(marks)
chart.draw_dots(x="hours", y="marks")
chart.draw_line_at_y(50, labelled="pass mark")
chart.shade_between_x(6, 8)
chart.draw_rectangle(x_between=[3, 5], y_between=[60, 80])
chart.draw_arrow(start=[2, 40], end=[4, 60])
chart.save("shapes.png")
\end{lstlisting}

\textbf{Left to matplotlib.} Polygons, paths and other patches. Hand the chart over and add them:

\begin{lstlisting}[style=py]
from matplotlib.patches import Polygon

marks = sdata.get_example("marks")
chart = sdata.chart(marks)
chart.draw_dots(x="hours", y="marks")
fig = chart.as_matplotlib()
fig.axes[0].add_patch(Polygon([[2, 40], [5, 90], [8, 60]], alpha=0.2, color="purple"))
fig.savefig("polygon.png")
\end{lstlisting}

\unsection{Charts in 3D}{from-matplotlib:3d}

\textbf{The job.} Draw a surface, dots or lines in three dimensions and choose the view. In the guide:
\href{https://matplotlib.org/stable/users/explain/toolkits/mplot3d.html}{The mplot3d toolkit}.

\textbf{In simplydata.} \code{draw\_surface} is \code{plot\_surface} (with \code{as\_wires=True} it is
\code{plot\_wireframe}), \code{draw\_dots\_3d} is \code{scatter} on 3D Axes, \code{draw\_lines\_3d}
is \code{plot} in 3D and \code{draw\_bars\_3d} is \code{bar3d}. \code{look\_from} is
\code{view\_init}, and \code{label\_z}, \code{set\_range\_z} and \code{put\_ticks\_z} work on the third
axis.

\begin{lstlisting}[style=py]
island = sdata.get_example("island")
chart = sdata.chart(island)
chart.draw_surface(x="east", y="north", z="height")
chart.look_from(height=40, around=30)
chart.label_z("Height (m)")
chart.save("surface.png")
\end{lstlisting}

\unsection{Moving charts}{from-matplotlib:animations}

\textbf{The job.} Make a chart that moves, one picture after another. In the guide:
\href{https://matplotlib.org/stable/users/explain/animations/animations.html}{Animations using
Matplotlib}.

\textbf{In simplydata.} \code{animate(over="year")} draws one picture for each value of a column:
matplotlib draws each picture and Pillow writes them into a .gif file.

\begin{lstlisting}[style=py]
growth = sdata.get_example("growth")
chart = sdata.chart(growth)
chart.draw_bars(x="country", y="population")
chart.animate(over="year")
chart.save("growth.gif")
\end{lstlisting}

\unsection{Saving, showing and web pages}{from-matplotlib:saving}

\textbf{The job.} Save the picture in a file type and size you need, show it in a window, or make a
chart you can point at in a web page. In the guide:
\href{https://matplotlib.org/stable/users/explain/figure/index.html}{Figures and backends}, with
\href{https://matplotlib.org/stable/users/explain/figure/backends.html}{Backends} and
\href{https://matplotlib.org/stable/users/explain/figure/interactive.html}{Interactive figures}.

\textbf{In simplydata.} \code{save} is \code{savefig}: the end of the file name picks the type
(.png, .pdf, .svg, .eps, .tif, .jpg, .gif or .html), the picture is 300 dots per inch and trimmed to
what is drawn, \code{resolution=} is \code{dpi}, and \code{see\_through=True} is
\code{transparent=True}. \code{set\_size} is \code{figsize}, in inches or cm. \code{show} is
\code{plt.show}. A .html file and \code{show\_in\_web\_page} make a web page, drawn by plotly from
what matplotlib drew; \code{show\_when\_pointing} chooses what it shows when the pointer rests on a
dot. In Jupyter a chart shows by itself.

\begin{lstlisting}[style=py]
marks = sdata.get_example("marks")
chart = sdata.chart(marks)
chart.draw_dots(x="hours", y="marks", colour="class")
chart.set_size(width=17, height=10, unit="cm")
chart.save("for_the_journal.pdf")
chart.save("for_a_slide.png", resolution=150, see_through=True)
chart.get_size(unit="cm")
\end{lstlisting}

% needs: plotly
\begin{lstlisting}[style=py]
marks = sdata.get_example("marks")
chart = sdata.chart(marks)
chart.draw_dots(x="hours", y="marks", colour="class")
chart.show_when_pointing("name", "subject")
chart.save("marks.html")
\end{lstlisting}

\textbf{The hand-over and web pages.} \code{chart.as\_matplotlib()} always gives the matplotlib
Figure, also for a chart you save as a web page: the web page is a copy that plotly draws from it.
Changes you make in matplotlib after the hand-over are not in \code{chart.save("marks.html")}.

\textbf{Left to matplotlib.} Choosing the backend, windows that react to the mouse or the keyboard
(events), and widgets such as sliders and buttons.

\unsection{Relationships: dots and lines}{from-matplotlib:relational}

\textbf{The job.} Show how one number goes with another: dots for each row, or a line through the
average at each x, with groups by colour, size or shape. In the tutorial:
\href{https://seaborn.pydata.org/tutorial/relational.html}{Visualizing statistical relationships}
(\code{relplot}, \code{scatterplot}, \code{lineplot}), and the overview of the
\href{https://seaborn.pydata.org/tutorial/function_overview.html}{plotting functions} and the
\href{https://seaborn.pydata.org/tutorial/data_structure.html}{data structures} seaborn takes.

\textbf{In simplydata.} \code{draw\_dots} is \code{scatterplot}: \code{colour=} is \code{hue=},
\code{size=} is \code{size=} and \code{shape=} is \code{style=}. \code{draw\_lines} is
\code{lineplot}; it draws the average of the rows at each x, and \code{error=} adds the band.
\code{relplot}'s panels are \code{draw\_one\_chart\_for\_each\_group\_in}.

\begin{lstlisting}[style=py]
marks = sdata.get_example("marks")
chart = sdata.chart(marks)
chart.draw_dots(x="hours", y="marks", colour="class", shape="subject")
chart.save("relationship.png")

school = sdata.get_example("school")
chart = sdata.chart(school)
chart.draw_lines(x="sleep", y="marks", colour="breakfast", error="confidence interval")
chart.save("average_line.png")
\end{lstlisting}

\unsection{Distributions}{from-matplotlib:distributions}

\textbf{The job.} Show how the values of a column are spread: a histogram, a smooth density, the
share of values up to each point, a rug, and two columns together. In the tutorial:
\href{https://seaborn.pydata.org/tutorial/distributions.html}{Visualizing distributions of data}
(\code{displot}, \code{histplot}, \code{kdeplot}, \code{ecdfplot}, \code{rugplot}).

\textbf{In simplydata.} \code{draw\_histogram} is \code{histplot} (\code{with\_bins=},
\code{show="density"}, \code{stacked=True}, \code{cumulative=True}); \code{draw\_density} is
\code{kdeplot}; \code{draw\_cumulative\_share} is \code{ecdfplot}; \code{draw\_rug\_lines} is
\code{rugplot}. With \code{y=} as well, \code{draw\_histogram} and \code{draw\_density} show two
columns together, and \code{draw\_side\_histograms} adds the edges of \code{jointplot}.
\code{fit\_distribution} draws a fitted curve over a histogram.

\begin{lstlisting}[style=py]
marks = sdata.get_example("marks")
chart = sdata.chart(marks)
chart.draw_histogram(x="marks", colour="class", show="density")
chart.draw_density(x="marks", colour="class")
chart.draw_rug_lines(x="marks")
chart.save("spread.png")

chart = sdata.chart(marks)
chart.draw_dots(x="hours", y="marks")
chart.draw_side_histograms(drawn_as="density")
chart.save("joint.png")
\end{lstlisting}

\unsection{Groups: categorical charts}{from-matplotlib:categorical}

\textbf{The job.} Compare the values of groups: every dot, boxes, violins, bars of an average with
an interval, points joined across groups, and counts. In the tutorial:
\href{https://seaborn.pydata.org/tutorial/categorical.html}{Visualizing categorical data}
(\code{catplot}, \code{stripplot}, \code{swarmplot}, \code{boxplot}, \code{violinplot},
\code{barplot}, \code{pointplot}, \code{countplot}).

\textbf{In simplydata.} \code{draw\_all\_dots} is \code{stripplot} and \code{swarmplot};
\code{draw\_boxes} is \code{boxplot}; \code{draw\_violins} is \code{violinplot} (\code{in\_halves=}
is \code{split=True}); \code{draw\_bars} is \code{barplot}, with \code{show=} for the average,
median, total or count; \code{draw\_interaction\_lines} is \code{pointplot}; \code{draw\_count\_bars}
is \code{countplot}. \code{catplot}'s panels are \code{draw\_one\_chart\_for\_each\_group\_in}.

\begin{lstlisting}[style=py]
marks = sdata.get_example("marks")
chart = sdata.chart(marks)
chart.draw_boxes(x="subject", y="marks")
chart.draw_all_dots(x="subject", y="marks")
chart.save("boxes_and_dots.png")

school = sdata.get_example("school")
chart = sdata.chart(school)
chart.draw_interaction_lines(x="class", y="marks", colour="breakfast")
chart.save("points.png")

chart = sdata.chart(school)
chart.draw_count_bars(x="class", colour="breakfast")
chart.save("counts.png")
\end{lstlisting}

\textbf{Left to seaborn.} \code{boxenplot}, the letter-value boxes for big tables. Hand an empty chart
over and let seaborn draw on its Axes:

\begin{lstlisting}[style=py]
import seaborn as sns

marks = sdata.get_example("marks")
chart = sdata.chart(marks)
chart.label_x("subject")
fig = chart.as_matplotlib()
sns.boxenplot(marks.as_pandas(), x="subject", y="marks", ax=fig.axes[0])
fig.savefig("boxen.png")
\end{lstlisting}

\unsection{Averages with error bars}{from-matplotlib:error-bars}

\textbf{The job.} Show how sure an average is (a confidence interval or a standard error) or how
spread the values are (a standard deviation). In the tutorial:
\href{https://seaborn.pydata.org/tutorial/error_bars.html}{Statistical estimation and error bars}.

\textbf{In simplydata.} \code{error=} on \code{draw\_bars}, \code{draw\_lines} and \code{draw\_dots}
takes \code{"confidence interval"}, \code{"standard error"} or \code{"standard deviation"}, as
seaborn's \code{errorbar=("ci", 95)}, \code{"se"} and \code{"sd"}. \code{draw\_error\_bars} adds them
to dots or lines. The interval is the t interval, not seaborn's bootstrap, so the same table always
gives the same bars, and a line under the chart says which one was drawn.

\begin{lstlisting}[style=py]
marks = sdata.get_example("marks")
chart = sdata.chart(marks)
chart.draw_bars(x="class", y="marks", show="average", error="confidence interval")
chart.save("error_bars.png")
\end{lstlisting}

\unsection{Lines that fit}{from-matplotlib:regression}

\textbf{The job.} Draw the line that fits the dots best, with its confidence band, a curve, a smooth
line, and what is left over. In the tutorial:
\href{https://seaborn.pydata.org/tutorial/regression.html}{Estimating regression fits}
(\code{regplot}, \code{lmplot}, \code{residplot}).

\textbf{In simplydata.} \code{fit\_line} is \code{regplot}; \code{show\_equation=True} writes the
equation and $R^2$ on the chart. \code{fit\_curve(of\_degree=2)} is \code{order=2}.
\code{draw\_smooth\_line} is \code{lowess=True}. \code{draw\_residual\_dots} is \code{residplot}.
For \code{lmplot}'s panels, add \code{draw\_one\_chart\_for\_each\_group\_in}.

\begin{lstlisting}[style=py]
marks = sdata.get_example("marks")
chart = sdata.chart(marks)
chart.draw_dots(x="hours", y="marks")
chart.fit_line(show_equation=True)
chart.save("fitted_line.png")

chart = sdata.chart(marks)
chart.draw_residual_dots(x="hours", y="marks")
chart.save("left_over.png")
\end{lstlisting}

\unsection{Grids of charts}{from-matplotlib:grids}

\textbf{The job.} One chart for each group in a grid, every column against every other, and a chart
with its edges. In the tutorial: \href{https://seaborn.pydata.org/tutorial/axis_grids.html}{Building
structured multi-plot grids} (\code{FacetGrid}, \code{PairGrid}, \code{pairplot}, \code{JointGrid},
\code{jointplot}).

\textbf{In simplydata.} \code{draw\_one\_chart\_for\_each\_group\_in} with \code{down=} and
\code{across=} is \code{FacetGrid} with \code{row=} and \code{col=}; \code{same\_y=False} is
\code{sharey=False}. \code{sdata.draw\_all\_pairs} is \code{pairplot}. \code{draw\_side\_histograms}
is \code{jointplot}.

\begin{lstlisting}[style=py]
marks = sdata.get_example("marks")
chart = sdata.chart(marks)
chart.draw_dots(x="hours", y="marks")
chart.draw_one_chart_for_each_group_in(down="class", across="subject")
chart.save("grid.png")

pairs = sdata.draw_all_pairs(marks, colour="class")
pairs.save("pairs.png")
\end{lstlisting}

\textbf{Left to seaborn.} Your own function mapped on each panel of a grid (\code{FacetGrid.map}), and
different kinds of charts above and below the diagonal of a \code{PairGrid}.

\unsection{The look of a chart}{from-matplotlib:aesthetics}

\textbf{The job.} Choose the overall look, take away lines of the frame, and scale everything for a
paper, a talk or a poster. In the tutorial:
\href{https://seaborn.pydata.org/tutorial/aesthetics.html}{Controlling figure aesthetics}
(\code{set\_theme}, \code{set\_style}, \code{despine}, \code{set\_context}).

\textbf{In simplydata.} The normal look (\code{"paper"}) already leaves out the top and right lines,
as \code{despine} does. \code{use\_style} is \code{set\_style} and \code{set\_context} in one word:
\code{"presentation"} is like seaborn's \code{"talk"}. \code{hide\_frame} takes the whole frame away,
\code{show\_grid} and \code{hide\_grid} the grid lines, and \code{use\_background\_colour} the
background.

\begin{lstlisting}[style=py]
weather = sdata.get_example("weather")
chart = sdata.chart(weather)
chart.draw_lines(x="month", y="rain", colour="city")
chart.use_style("dark")
chart.show_grid()
chart.save("dark.png")
\end{lstlisting}

\unsection{Colour palettes}{from-matplotlib:palettes}

\textbf{The job.} Choose colours for groups (a qualitative palette), for numbers that go from low to
high (sequential) and for numbers on both sides of a middle (diverging). In the tutorial:
\href{https://seaborn.pydata.org/tutorial/color_palettes.html}{Choosing color palettes}.

\textbf{In simplydata.} The normal colours for groups are the Okabe-Ito set, safe for readers who do
not see some colours. \code{use\_colours} is \code{color\_palette} and \code{set\_palette}: a named set
(\code{"Set2"}, \code{"deep"}, \code{"colorblind"}), your own colours, or \code{how\_many=} colours
from a colour map. \code{use\_colour\_scale} is the \code{cmap} of a heat map; \code{"blue to red"}
with \code{middle=0} is a diverging palette centred on 0.

\begin{lstlisting}[style=py]
marks = sdata.get_example("marks")
chart = sdata.chart(marks)
chart.draw_bars(x="class", y="marks", colour="subject")
chart.use_colours("viridis", how_many=3)
chart.save("palette.png")
\end{lstlisting}

\unsection{The jobs left to matplotlib and seaborn}{from-matplotlib:left-out}

simplydata draws the charts most people need in research and teaching, with one word each. These jobs
are left out on purpose. Each one is a hand-over: \code{fig = chart.as\_matplotlib()}, then
matplotlib or seaborn code on \code{fig.axes}, then \code{fig.savefig(...)}.
\begin{itemize}
  \item \textbf{Arrows and vector fields} (\code{quiver}, \code{barbs}, \code{streamplot}).
  \item \textbf{Polygons, paths and other patches} (\code{fill}, \code{Polygon}), shown above.
  \item \textbf{A table drawn inside a chart} (\code{ax.table}).
  \item \textbf{Text of your own shape}: arrow styles, text along a path, many box styles (shown above).
  \item \textbf{Tick formatters as your own function} and scales of your own (shown above).
  \item \textbf{Axes placed by hand} (\code{fig.add\_axes}) and margins (\code{ax.margins}).
  \item \textbf{Images} (\code{imread}, \code{imsave}, \code{figimage}); images get their own words in
    a later phase.
  \item \textbf{Interaction}: backends, mouse and key events, widgets (\code{Slider}, \code{Button}).
  \item \textbf{More 3D}: voxels, arrows and contour sets in 3D.
  \item \textbf{seaborn's letter-value boxes} (\code{boxenplot}, shown above), a function of your own on
    each panel of a grid (\code{FacetGrid.map}) and the objects interface (\code{seaborn.objects}).
  \item \textbf{Any \code{rcParams} setting} beyond the looks, fonts and sizes of \code{use\_style},
    \code{use\_font} and \code{use\_font\_size} (shown above).
\end{itemize}

One more example: a vector field over a contour map.

\begin{lstlisting}[style=py]
island = sdata.get_example("island")
chart = sdata.chart(island)
chart.draw_contour(x="east", y="north", z="height")
fig = chart.as_matplotlib()
ax = fig.axes[0]
ax.quiver([2000, 4000], [1000, 3000], [1, -1], [1, 0.5], color="white")
fig.savefig("vectors.png")
\end{lstlisting}
).
% The topics follow the order of the matplotlib Quick start guide and User guide, then the seaborn tutorial, used as a
% checklist only: the text and the examples are our own, and each section links to the official page. Every example in
% a py box is run by tests/test_from_matplotlib_chapter.py (with "import simplydata as sdata" in front, and matplotlib's
% Agg backend), so it cannot go out of date. A line "% needs: plotly" just before a box says the box needs that package.

\unchapter{From matplotlib and seaborn to simplydata}{from-matplotlib}

This chapter is for you if you already know matplotlib or seaborn, or if you want to learn them after
simplydata. It walks through the topics of the official guides, in their order: the
\href{https://matplotlib.org/stable/users/explain/quick_start.html}{matplotlib Quick start guide} and
\href{https://matplotlib.org/stable/users/index.html}{User guide}, then the
\href{https://seaborn.pydata.org/tutorial.html}{seaborn tutorial}. For each topic it says what the job
is, which simplydata words do it, and a short example on the practice tables. Some jobs are left to
matplotlib and seaborn on purpose. For those, the chapter shows how to hand the chart over and keep
going in matplotlib.

Every example starts with \code{import simplydata as sdata}. It is written once here, not in each
box. The examples use the practice tables of \code{sdata.get\_example}, so you can run them as they
are.

\unsection{Five things to know first}{from-matplotlib:first}

\begin{enumerate}
  \item \textbf{A chart is a list of steps.} In matplotlib each call draws at once. In simplydata,
    \code{chart.draw\_dots(...)} and \code{chart.label\_x(...)} only write down a step. The chart is
    drawn when you save it, show it, or look at it in Jupyter. So the order of the steps rarely
    matters, and \code{chart.undo()} takes the last step back.
  \item \textbf{There is no figure and no axes to hold.} simplydata makes the matplotlib Figure and
    its Axes for you each time it draws. You never write \code{fig, ax = plt.subplots()}.
  \item \textbf{Columns by name.} You give the table once, to \code{sdata.chart(table)}, and then the
    column names, as seaborn does with \code{data=}. A word works on a pandas or a Polars table.
  \item \textbf{One word for one job.} matplotlib often has several ways to do the same thing
    (\code{plt.title}, \code{ax.set\_title}, \code{fig.suptitle}). simplydata has one word,
    \code{write\_title\_on\_top}, and \code{sdata.help("ax.set\_title")} names it.
  \item \textbf{The door is always open.} \code{fig = chart.as\_matplotlib()} gives the chart as a
    plain matplotlib Figure. Its charts are the Axes in \code{fig.axes}, so \code{ax = fig.axes[0]}.
    Then any matplotlib code works. It is a copy: changing \code{fig} does not change the chart.
    Save it with \code{fig.savefig(...)}.
\end{enumerate}

\begin{lstlisting}[style=py]
marks = sdata.get_example("marks")
chart = sdata.chart(marks)
chart.draw_dots(x="hours", y="marks")
fig = chart.as_matplotlib()              # a matplotlib Figure: use any matplotlib code here
ax = fig.axes[0]
ax.set_facecolor("whitesmoke")
fig.savefig("handed_over.png", dpi=150)
\end{lstlisting}

Do you know a matplotlib or seaborn name, but not the simplydata word? Ask:
\code{sdata.help("ax.set\_xlabel")}, \code{sdata.help("plt.hist")} or
\code{sdata.help("sns.lineplot")}. The name can have \code{plt.}, \code{ax.}, \code{fig.} or
\code{sns.} in front. The full list is in the section ``Coming from another tool?'' of the simplydata
dictionary
(\href{https://github.com/bioforkgithub/simplydata/blob/main/docs/manual/simplydataDictionary.pdf}%
{simplydataDictionary.pdf}).

\unsection{A first chart}{from-matplotlib:first-chart}

\textbf{The job.} Draw a line from two columns, give it a title and axis names, and save it as a
picture. In the guide: \href{https://matplotlib.org/stable/users/explain/quick_start.html}{Quick
start guide}, the simple example at its start.

\textbf{In simplydata.} \code{draw\_lines} is \code{ax.plot}; \code{write\_title\_on\_top},
\code{label\_x} and \code{label\_y} are \code{set\_title}, \code{set\_xlabel} and
\code{set\_ylabel}; \code{save} is \code{fig.savefig}. Without \code{label\_x}, the column name is the
axis name.

\begin{lstlisting}[style=py]
weather = sdata.get_example("weather")
chart = sdata.chart(weather)
chart.draw_lines(x="month", y="temperature", colour="city")
chart.write_title_on_top("Temperature through the year")
chart.label_y("Temperature (degrees C)")
chart.save("first_chart.png")
\end{lstlisting}

\unsection{The parts of a figure}{from-matplotlib:parts}

\textbf{The job.} Know the parts matplotlib draws: the Figure (the whole picture), the Axes (one
chart inside it), each Axis (the x and y scales) and the Artists (every line, dot, bar and piece of
text). In the guide: \href{https://matplotlib.org/stable/users/explain/quick_start.html}{Quick start
guide}, ``Parts of a Figure'', and
\href{https://matplotlib.org/stable/users/explain/artists/index.html}{Artists}.

\textbf{In simplydata.} A simplydata chart is one Figure with one Axes, or several Axes after
\code{draw\_one\_chart\_for\_each\_group\_in}. You do not touch the parts. You can read them back as
tables: \code{get\_numbers\_drawn} reads the lines, dots and bars, \code{get\_text\_drawn} the text,
and \code{get\_scales\_drawn} each axis. \code{as\_matplotlib} gives the parts themselves.

\begin{lstlisting}[style=py]
marks = sdata.get_example("marks")
chart = sdata.chart(marks)
chart.draw_bars(x="class", y="marks")
chart.get_numbers_drawn()                # one row for each bar: where it is and how high
chart.get_scales_drawn()                 # the range and kind of each axis
fig = chart.as_matplotlib()
print(len(fig.axes), "Axes in the Figure")
\end{lstlisting}

\unsection{The data a chart takes}{from-matplotlib:data}

\textbf{The job.} Give the numbers to draw. matplotlib takes lists or NumPy arrays, or column names
with \code{data=}. In the guide: \href{https://matplotlib.org/stable/users/explain/quick_start.html}%
{Quick start guide}, ``Types of inputs to plotting functions''.

\textbf{In simplydata.} A table and column names is the usual way: \code{sdata.chart(table)} reads a
table, a file name or a pandas or Polars DataFrame. Without a table, give lists to the step itself.
Numbers made with \code{sdata.number\_list} start their own chart.

\begin{lstlisting}[style=py]
chart = sdata.chart()                    # no table: lists go to the step
chart.draw_lines(x=[1, 2, 3, 4], y=[10, 12, 9, 14])
chart.save("from_lists.png")

numbers = sdata.number_list([4, 8, 15, 16, 23, 42])
chart = numbers.start_chart()
chart.draw_histogram(x="value")
chart.save("from_numbers.png")
\end{lstlisting}

\unsection{Two ways of writing, one in simplydata}{from-matplotlib:styles}

\textbf{The job.} matplotlib can be written in two styles: the explicit one (\code{fig, ax =
plt.subplots()}, then \code{ax.plot(...)}) and the implicit one with \code{pyplot}
(\code{plt.plot(...)}, which draws on the current Axes). In the guide:
\href{https://matplotlib.org/stable/users/explain/quick_start.html}{Quick start guide}, ``Coding
styles''.

\textbf{In simplydata.} There is one style. The chart is the object you name, as in the explicit style,
and nothing is ``current''. So two charts never draw on each other. To make many charts the same way,
write a small function that gives back a chart, as the guide suggests for matplotlib.

\begin{lstlisting}[style=py]
def marks_chart(subject):
    marks = sdata.get_example("marks")
    marks.keep_rows_where("subject", is_equal_to=subject)
    chart = sdata.chart(marks)
    chart.draw_dots(x="hours", y="marks")
    chart.write_title_on_top(subject)
    return chart


marks_chart("Maths").save("maths.png")
marks_chart("Science").save("science.png")
\end{lstlisting}

\unsection{Colours, lines and dots}{from-matplotlib:styling}

\textbf{The job.} Choose the colour, width and dash of a line, and the size, shape and transparency of
dots. In the guide: \href{https://matplotlib.org/stable/users/explain/quick_start.html}{Quick start
guide}, ``Styling Artists'', and
\href{https://matplotlib.org/stable/users/explain/colors/colors.html}{Specifying colors}.

\textbf{In simplydata.} Options with plain names on the draw words: \code{thickness=} is
\code{linewidth}, \code{dashed=True} and \code{dotted=True} are \code{linestyle}, \code{size=} is the
dot size (a column, for sizes that change), \code{transparency=} is one minus \code{alpha}.
\code{use\_colour} sets one colour, \code{use\_colours} one colour for each group, and
\code{use\_shapes} a shape for each group. The groups come from \code{colour="column"}, as seaborn's
\code{hue=}.

\begin{lstlisting}[style=py]
weather = sdata.get_example("weather")
chart = sdata.chart(weather)
chart.draw_lines(x="month", y="temperature", colour="city", thickness=3, dashed="city")
chart.use_colours("black", "darkorange")
chart.save("styled_lines.png")

marks = sdata.get_example("marks")
chart = sdata.chart(marks)
chart.draw_dots(x="hours", y="marks", colour="class", size="height", transparency=0.3)
chart.use_shapes("circle", "square", "triangle")
chart.save("styled_dots.png")
\end{lstlisting}

\unsection{Titles, labels, notes and legends}{from-matplotlib:labels}

\textbf{The job.} Name the axes, write text on the chart, point at a place with an arrow, and place
the legend. In the guide: \href{https://matplotlib.org/stable/users/explain/quick_start.html}{Quick
start guide}, ``Labelling plots'';
\href{https://matplotlib.org/stable/users/explain/text/text_intro.html}{Text in Matplotlib},
\href{https://matplotlib.org/stable/users/explain/text/annotations.html}{Annotations} and the
\href{https://matplotlib.org/stable/users/explain/axes/legend_guide.html}{Legend guide}.

\textbf{In simplydata.} \code{label\_x}, \code{label\_y}, \code{write\_title\_on\_top} and
\code{write\_caption\_at\_bottom} write the text around the chart. \code{write\_note} is
\code{ax.text}, and with \code{with\_arrow=True} it is \code{ax.annotate}. The text goes to
matplotlib as you write it, so matplotlib's maths between dollar signs works too. A legend appears by
itself when \code{colour=} makes groups; \code{put\_legend\_inside}, \code{put\_legend\_at\_bottom}
and the other \code{put\_legend\_} words place it, \code{label\_legend} names it,
\code{order\_legend} orders it and \code{hide\_legend} hides it.

\begin{lstlisting}[style=py]
growth = sdata.get_example("growth")
chart = sdata.chart(growth)
chart.draw_lines(x="year", y="population", colour="country")
chart.label_y("People (millions)")
chart.label_x("Year, $t$")
chart.write_note("fastest growth", x=2020, y=12, with_arrow=True)
chart.write_caption_at_bottom("Made-up data")
chart.put_legend_inside("top left")
chart.label_legend("Country")
chart.save("labels.png")
\end{lstlisting}

\textbf{Left to matplotlib.} Arrows with your own shape and connection, text along a path, and
boxes of many styles. Hand the chart over:

\begin{lstlisting}[style=py]
growth = sdata.get_example("growth")
chart = sdata.chart(growth)
chart.draw_lines(x="year", y="population", colour="country")
fig = chart.as_matplotlib()
ax = fig.axes[0]
ax.annotate("a bend", xy=(2010, 6), xytext=(2003, 10),
            arrowprops={"arrowstyle": "-|>", "connectionstyle": "arc3,rad=0.3"},
            bbox={"boxstyle": "round", "facecolor": "lightyellow"})
fig.savefig("own_arrow.png")
\end{lstlisting}

\unsection{Scales, ticks, dates and groups on an axis}{from-matplotlib:axes}

\textbf{The job.} Choose the range and the scale of an axis (log, square root), where the ticks go and
how their numbers are written, dates and text on an axis, and a second axis. In the guide:
\href{https://matplotlib.org/stable/users/explain/quick_start.html}{Quick start guide}, ``Axis
scales and ticks'';
\href{https://matplotlib.org/stable/users/explain/axes/axes_scales.html}{Axis scales} and
\href{https://matplotlib.org/stable/users/explain/axes/axes_ticks.html}{Axis ticks}.

\textbf{In simplydata.} \code{set\_range\_x} and \code{set\_range\_y} are \code{set\_xlim} and
\code{set\_ylim}; \code{use\_log\_scale\_y} and \code{use\_scale\_y("square root")} are
\code{set\_yscale}; \code{put\_ticks\_x} is \code{set\_xticks} (with \code{labels=} it is
\code{set\_xticklabels}, and \code{every=} a \code{MultipleLocator}); \code{write\_ticks\_y\_as}
chooses a formatter; \code{rotate\_labels\_x} turns the labels; \code{flip\_y} is
\code{invert\_yaxis}. Dates and text on an axis need nothing: simplydata finds them. Months and
weekdays come in calendar order, and \code{order\_x} sets any other order. \code{add\_second\_y} is
\code{twinx}, and \code{add\_second\_y(formula=...)} is \code{secondary\_yaxis}.

\begin{lstlisting}[style=py]
shop = sdata.get_example("shop")
chart = sdata.chart(shop)
chart.draw_lines(x="date", y="sales", colour="shop")
chart.write_ticks_x_as("Jan 2026")
chart.write_ticks_y_as("thousands")
chart.rotate_labels_x(45)
chart.save("dates.png")

marks = sdata.get_example("marks")
chart = sdata.chart(marks)
chart.draw_bars(x="class", y="marks")
chart.order_x("C", "B", "A")
chart.set_range_y(0, 100)
chart.put_ticks_y(0, 25, 50, 75, 100)
chart.save("groups.png")

weather = sdata.get_example("weather")
chart = sdata.chart(weather)
chart.draw_lines(x="month", y="temperature", colour="city")
chart.add_second_y(formula="y * 9 / 5 + 32")
chart.save("second_axis.png")
\end{lstlisting}

\textbf{Left to matplotlib.} A tick formatter written as your own function, and scales of your own
(\code{FuncScale}). Hand the chart over:

\begin{lstlisting}[style=py]
from matplotlib.ticker import FuncFormatter

marks = sdata.get_example("marks")
chart = sdata.chart(marks)
chart.draw_dots(x="hours", y="marks")
fig = chart.as_matplotlib()
fig.axes[0].xaxis.set_major_formatter(FuncFormatter(lambda value, place: f"{value:.0f} h"))
fig.savefig("own_ticks.png")
\end{lstlisting}

\unsection{Numbers as colours}{from-matplotlib:colour-maps}

\textbf{The job.} Show a number as a colour (a heat map or a contour map), choose the colour map and
how numbers map to colours, and add a colour bar. In the guide:
\href{https://matplotlib.org/stable/users/explain/quick_start.html}{Quick start guide}, ``Color
mapped data''; \href{https://matplotlib.org/stable/users/explain/colors/colormaps.html}{Choosing
Colormaps} and \href{https://matplotlib.org/stable/users/explain/colors/colormapnorms.html}{Colormap
normalization}.

\textbf{In simplydata.} \code{draw\_heatmap} is \code{imshow} or \code{pcolormesh} (seaborn
\code{heatmap}), and \code{draw\_contour} is \code{contourf}, also for scattered points.
\code{use\_colour\_scale} is \code{cmap=}: it chooses the colour map. Its options are the
normalizations: \code{middle=0} puts the middle colour at 0 (a \code{TwoSlopeNorm}), \code{log=True}
spreads the colours on a log scale, \code{steps=} cuts the scale into bands, and \code{between=} sets
\code{vmin} and \code{vmax}. The colour bar comes by itself; \code{label\_colour\_bar} names it and
\code{hide\_colour\_bar} hides it.

\begin{lstlisting}[style=py]
weather = sdata.get_example("weather")
chart = sdata.chart(weather)
chart.draw_heatmap(x="month", y="city", value="rain", with_numbers=True)
chart.use_colour_scale("blues", between=[0, 300])
chart.label_colour_bar("Rain (mm)")
chart.save("heatmap.png")

island = sdata.get_example("island")
chart = sdata.chart(island)
chart.draw_contour(x="east", y="north", z="height")
chart.use_colour_scale("terrain", steps=8)
chart.save("contour.png")
\end{lstlisting}

\unsection{Several charts in one picture}{from-matplotlib:several}

\textbf{The job.} Put several Axes in one Figure: a grid of charts, one chart for each group, charts of
different sizes, or a small chart inside a big one. In the guide:
\href{https://matplotlib.org/stable/users/explain/quick_start.html}{Quick start guide}, ``Working
with multiple Figures and Axes'';
\href{https://matplotlib.org/stable/users/explain/axes/arranging_axes.html}{Arranging multiple Axes in
a Figure} and the
\href{https://matplotlib.org/stable/users/explain/axes/constrainedlayout_guide.html}{Constrained layout
guide}.

\textbf{In simplydata.} \code{draw\_one\_chart\_for\_each\_group\_in} is seaborn's
\code{FacetGrid} (\code{col=}, \code{row=}, \code{col\_wrap=}). \code{sdata.put\_side\_by\_side},
\code{sdata.put\_one\_above\_other} and \code{sdata.put\_in\_rows\_and\_columns} are
\code{plt.subplots} and \code{subplot\_mosaic} for charts you made one by one; a nested list makes
a chart span several places, and \code{with\_letters=True} writes A, B, C. \code{put\_chart\_inside} is
\code{inset\_axes}, and \code{add\_zoom\_box} is \code{indicate\_inset\_zoom}. The layout is always
matplotlib's constrained layout, so labels do not overlap.

\begin{lstlisting}[style=py]
marks = sdata.get_example("marks")
chart = sdata.chart(marks)
chart.draw_dots(x="hours", y="marks")
chart.draw_one_chart_for_each_group_in("subject", per_row=3)
chart.save("one_per_group.png")

dots = sdata.chart(marks).draw_dots(x="hours", y="marks")
bars = sdata.chart(marks).draw_bars(x="class", y="marks")
boxes = sdata.chart(marks).draw_boxes(x="subject", y="marks")
page = sdata.put_in_rows_and_columns([[dots, dots], [bars, boxes]], with_letters=True)
page.save("page.png")
\end{lstlisting}

\textbf{Left to matplotlib.} Axes placed by hand at any place in the Figure (\code{fig.add\_axes}).
\code{page.as\_matplotlib()} gives the page as one Figure: each chart is a SubFigure in
\code{fig.subfigs}, and \code{fig.axes} holds every Axes of every chart.

\unsection{Looks and settings}{from-matplotlib:style}

\textbf{The job.} Change the look of every chart at once: a style sheet, the fonts and the sizes of the
text (matplotlib's \code{rcParams}). In the guide:
\href{https://matplotlib.org/stable/users/explain/customizing.html}{Customizing Matplotlib with style
sheets and rcParams}.

\textbf{In simplydata.} \code{use\_style} picks one of five looks: \code{"paper"} (normal),
\code{"school"}, \code{"presentation"}, \code{"dark"} and \code{"black and white"}.
\code{use\_font} and \code{use\_font\_size} set the text. Each chart keeps its own look; nothing
changes matplotlib's settings for other code.

\begin{lstlisting}[style=py]
marks = sdata.get_example("marks")
chart = sdata.chart(marks)
chart.draw_boxes(x="class", y="marks")
chart.use_style("presentation")
chart.use_font("serif")
chart.use_font_size(16, title=22)
chart.save("presentation.png")
\end{lstlisting}

\textbf{Left to matplotlib.} Any other setting of \code{rcParams}. A setting such as the colour of the
axis lines is used while matplotlib draws, so change the drawn parts after the hand-over:

\begin{lstlisting}[style=py]
marks = sdata.get_example("marks")
chart = sdata.chart(marks)
chart.draw_boxes(x="class", y="marks")
fig = chart.as_matplotlib()
for side in fig.axes[0].spines.values():
    side.set_color("steelblue")
fig.savefig("own_settings.png")
\end{lstlisting}

\unsection{Lines, bands and shapes on a chart}{from-matplotlib:marks}

\textbf{The job.} Draw a reference line, shade a range, draw a rectangle, a circle or an arrow. In the
guide: \href{https://matplotlib.org/stable/plot_types/index.html}{Plot types} and
\href{https://matplotlib.org/stable/users/explain/text/annotations.html}{Annotations}.

\textbf{In simplydata.} \code{draw\_line\_at\_y} is \code{axhline} (with \code{between=} it is
\code{hlines}), \code{draw\_line\_at\_x} is \code{axvline}, \code{draw\_line\_with\_slope} is
\code{axline}, \code{shade\_between\_x} is \code{axvspan}, \code{draw\_band} is
\code{fill\_between}, \code{draw\_rectangle} and \code{draw\_circle} are the \code{Rectangle} and
\code{Circle} patches, and \code{draw\_arrow} is an arrow from one point to another.

\begin{lstlisting}[style=py]
marks = sdata.get_example("marks")
chart = sdata.chart(marks)
chart.draw_dots(x="hours", y="marks")
chart.draw_line_at_y(50, labelled="pass mark")
chart.shade_between_x(6, 8)
chart.draw_rectangle(x_between=[3, 5], y_between=[60, 80])
chart.draw_arrow(start=[2, 40], end=[4, 60])
chart.save("shapes.png")
\end{lstlisting}

\textbf{Left to matplotlib.} Polygons, paths and other patches. Hand the chart over and add them:

\begin{lstlisting}[style=py]
from matplotlib.patches import Polygon

marks = sdata.get_example("marks")
chart = sdata.chart(marks)
chart.draw_dots(x="hours", y="marks")
fig = chart.as_matplotlib()
fig.axes[0].add_patch(Polygon([[2, 40], [5, 90], [8, 60]], alpha=0.2, color="purple"))
fig.savefig("polygon.png")
\end{lstlisting}

\unsection{Charts in 3D}{from-matplotlib:3d}

\textbf{The job.} Draw a surface, dots or lines in three dimensions and choose the view. In the guide:
\href{https://matplotlib.org/stable/users/explain/toolkits/mplot3d.html}{The mplot3d toolkit}.

\textbf{In simplydata.} \code{draw\_surface} is \code{plot\_surface} (with \code{as\_wires=True} it is
\code{plot\_wireframe}), \code{draw\_dots\_3d} is \code{scatter} on 3D Axes, \code{draw\_lines\_3d}
is \code{plot} in 3D and \code{draw\_bars\_3d} is \code{bar3d}. \code{look\_from} is
\code{view\_init}, and \code{label\_z}, \code{set\_range\_z} and \code{put\_ticks\_z} work on the third
axis.

\begin{lstlisting}[style=py]
island = sdata.get_example("island")
chart = sdata.chart(island)
chart.draw_surface(x="east", y="north", z="height")
chart.look_from(height=40, around=30)
chart.label_z("Height (m)")
chart.save("surface.png")
\end{lstlisting}

\unsection{Moving charts}{from-matplotlib:animations}

\textbf{The job.} Make a chart that moves, one picture after another. In the guide:
\href{https://matplotlib.org/stable/users/explain/animations/animations.html}{Animations using
Matplotlib}.

\textbf{In simplydata.} \code{animate(over="year")} draws one picture for each value of a column:
matplotlib draws each picture and Pillow writes them into a .gif file.

\begin{lstlisting}[style=py]
growth = sdata.get_example("growth")
chart = sdata.chart(growth)
chart.draw_bars(x="country", y="population")
chart.animate(over="year")
chart.save("growth.gif")
\end{lstlisting}

\unsection{Saving, showing and web pages}{from-matplotlib:saving}

\textbf{The job.} Save the picture in a file type and size you need, show it in a window, or make a
chart you can point at in a web page. In the guide:
\href{https://matplotlib.org/stable/users/explain/figure/index.html}{Figures and backends}, with
\href{https://matplotlib.org/stable/users/explain/figure/backends.html}{Backends} and
\href{https://matplotlib.org/stable/users/explain/figure/interactive.html}{Interactive figures}.

\textbf{In simplydata.} \code{save} is \code{savefig}: the end of the file name picks the type
(.png, .pdf, .svg, .eps, .tif, .jpg, .gif or .html), the picture is 300 dots per inch and trimmed to
what is drawn, \code{resolution=} is \code{dpi}, and \code{see\_through=True} is
\code{transparent=True}. \code{set\_size} is \code{figsize}, in inches or cm. \code{show} is
\code{plt.show}. A .html file and \code{show\_in\_web\_page} make a web page, drawn by plotly from
what matplotlib drew; \code{show\_when\_pointing} chooses what it shows when the pointer rests on a
dot. In Jupyter a chart shows by itself.

\begin{lstlisting}[style=py]
marks = sdata.get_example("marks")
chart = sdata.chart(marks)
chart.draw_dots(x="hours", y="marks", colour="class")
chart.set_size(width=17, height=10, unit="cm")
chart.save("for_the_journal.pdf")
chart.save("for_a_slide.png", resolution=150, see_through=True)
chart.get_size(unit="cm")
\end{lstlisting}

% needs: plotly
\begin{lstlisting}[style=py]
marks = sdata.get_example("marks")
chart = sdata.chart(marks)
chart.draw_dots(x="hours", y="marks", colour="class")
chart.show_when_pointing("name", "subject")
chart.save("marks.html")
\end{lstlisting}

\textbf{The hand-over and web pages.} \code{chart.as\_matplotlib()} always gives the matplotlib
Figure, also for a chart you save as a web page: the web page is a copy that plotly draws from it.
Changes you make in matplotlib after the hand-over are not in \code{chart.save("marks.html")}.

\textbf{Left to matplotlib.} Choosing the backend, windows that react to the mouse or the keyboard
(events), and widgets such as sliders and buttons.

\unsection{Relationships: dots and lines}{from-matplotlib:relational}

\textbf{The job.} Show how one number goes with another: dots for each row, or a line through the
average at each x, with groups by colour, size or shape. In the tutorial:
\href{https://seaborn.pydata.org/tutorial/relational.html}{Visualizing statistical relationships}
(\code{relplot}, \code{scatterplot}, \code{lineplot}), and the overview of the
\href{https://seaborn.pydata.org/tutorial/function_overview.html}{plotting functions} and the
\href{https://seaborn.pydata.org/tutorial/data_structure.html}{data structures} seaborn takes.

\textbf{In simplydata.} \code{draw\_dots} is \code{scatterplot}: \code{colour=} is \code{hue=},
\code{size=} is \code{size=} and \code{shape=} is \code{style=}. \code{draw\_lines} is
\code{lineplot}; it draws the average of the rows at each x, and \code{error=} adds the band.
\code{relplot}'s panels are \code{draw\_one\_chart\_for\_each\_group\_in}.

\begin{lstlisting}[style=py]
marks = sdata.get_example("marks")
chart = sdata.chart(marks)
chart.draw_dots(x="hours", y="marks", colour="class", shape="subject")
chart.save("relationship.png")

school = sdata.get_example("school")
chart = sdata.chart(school)
chart.draw_lines(x="sleep", y="marks", colour="breakfast", error="confidence interval")
chart.save("average_line.png")
\end{lstlisting}

\unsection{Distributions}{from-matplotlib:distributions}

\textbf{The job.} Show how the values of a column are spread: a histogram, a smooth density, the
share of values up to each point, a rug, and two columns together. In the tutorial:
\href{https://seaborn.pydata.org/tutorial/distributions.html}{Visualizing distributions of data}
(\code{displot}, \code{histplot}, \code{kdeplot}, \code{ecdfplot}, \code{rugplot}).

\textbf{In simplydata.} \code{draw\_histogram} is \code{histplot} (\code{with\_bins=},
\code{show="density"}, \code{stacked=True}, \code{cumulative=True}); \code{draw\_density} is
\code{kdeplot}; \code{draw\_cumulative\_share} is \code{ecdfplot}; \code{draw\_rug\_lines} is
\code{rugplot}. With \code{y=} as well, \code{draw\_histogram} and \code{draw\_density} show two
columns together, and \code{draw\_side\_histograms} adds the edges of \code{jointplot}.
\code{fit\_distribution} draws a fitted curve over a histogram.

\begin{lstlisting}[style=py]
marks = sdata.get_example("marks")
chart = sdata.chart(marks)
chart.draw_histogram(x="marks", colour="class", show="density")
chart.draw_density(x="marks", colour="class")
chart.draw_rug_lines(x="marks")
chart.save("spread.png")

chart = sdata.chart(marks)
chart.draw_dots(x="hours", y="marks")
chart.draw_side_histograms(drawn_as="density")
chart.save("joint.png")
\end{lstlisting}

\unsection{Groups: categorical charts}{from-matplotlib:categorical}

\textbf{The job.} Compare the values of groups: every dot, boxes, violins, bars of an average with
an interval, points joined across groups, and counts. In the tutorial:
\href{https://seaborn.pydata.org/tutorial/categorical.html}{Visualizing categorical data}
(\code{catplot}, \code{stripplot}, \code{swarmplot}, \code{boxplot}, \code{violinplot},
\code{barplot}, \code{pointplot}, \code{countplot}).

\textbf{In simplydata.} \code{draw\_all\_dots} is \code{stripplot} and \code{swarmplot};
\code{draw\_boxes} is \code{boxplot}; \code{draw\_violins} is \code{violinplot} (\code{in\_halves=}
is \code{split=True}); \code{draw\_bars} is \code{barplot}, with \code{show=} for the average,
median, total or count; \code{draw\_interaction\_lines} is \code{pointplot}; \code{draw\_count\_bars}
is \code{countplot}. \code{catplot}'s panels are \code{draw\_one\_chart\_for\_each\_group\_in}.

\begin{lstlisting}[style=py]
marks = sdata.get_example("marks")
chart = sdata.chart(marks)
chart.draw_boxes(x="subject", y="marks")
chart.draw_all_dots(x="subject", y="marks")
chart.save("boxes_and_dots.png")

school = sdata.get_example("school")
chart = sdata.chart(school)
chart.draw_interaction_lines(x="class", y="marks", colour="breakfast")
chart.save("points.png")

chart = sdata.chart(school)
chart.draw_count_bars(x="class", colour="breakfast")
chart.save("counts.png")
\end{lstlisting}

\textbf{Left to seaborn.} \code{boxenplot}, the letter-value boxes for big tables. Hand an empty chart
over and let seaborn draw on its Axes:

\begin{lstlisting}[style=py]
import seaborn as sns

marks = sdata.get_example("marks")
chart = sdata.chart(marks)
chart.label_x("subject")
fig = chart.as_matplotlib()
sns.boxenplot(marks.as_pandas(), x="subject", y="marks", ax=fig.axes[0])
fig.savefig("boxen.png")
\end{lstlisting}

\unsection{Averages with error bars}{from-matplotlib:error-bars}

\textbf{The job.} Show how sure an average is (a confidence interval or a standard error) or how
spread the values are (a standard deviation). In the tutorial:
\href{https://seaborn.pydata.org/tutorial/error_bars.html}{Statistical estimation and error bars}.

\textbf{In simplydata.} \code{error=} on \code{draw\_bars}, \code{draw\_lines} and \code{draw\_dots}
takes \code{"confidence interval"}, \code{"standard error"} or \code{"standard deviation"}, as
seaborn's \code{errorbar=("ci", 95)}, \code{"se"} and \code{"sd"}. \code{draw\_error\_bars} adds them
to dots or lines. The interval is the t interval, not seaborn's bootstrap, so the same table always
gives the same bars, and a line under the chart says which one was drawn.

\begin{lstlisting}[style=py]
marks = sdata.get_example("marks")
chart = sdata.chart(marks)
chart.draw_bars(x="class", y="marks", show="average", error="confidence interval")
chart.save("error_bars.png")
\end{lstlisting}

\unsection{Lines that fit}{from-matplotlib:regression}

\textbf{The job.} Draw the line that fits the dots best, with its confidence band, a curve, a smooth
line, and what is left over. In the tutorial:
\href{https://seaborn.pydata.org/tutorial/regression.html}{Estimating regression fits}
(\code{regplot}, \code{lmplot}, \code{residplot}).

\textbf{In simplydata.} \code{fit\_line} is \code{regplot}; \code{show\_equation=True} writes the
equation and $R^2$ on the chart. \code{fit\_curve(of\_degree=2)} is \code{order=2}.
\code{draw\_smooth\_line} is \code{lowess=True}. \code{draw\_residual\_dots} is \code{residplot}.
For \code{lmplot}'s panels, add \code{draw\_one\_chart\_for\_each\_group\_in}.

\begin{lstlisting}[style=py]
marks = sdata.get_example("marks")
chart = sdata.chart(marks)
chart.draw_dots(x="hours", y="marks")
chart.fit_line(show_equation=True)
chart.save("fitted_line.png")

chart = sdata.chart(marks)
chart.draw_residual_dots(x="hours", y="marks")
chart.save("left_over.png")
\end{lstlisting}

\unsection{Grids of charts}{from-matplotlib:grids}

\textbf{The job.} One chart for each group in a grid, every column against every other, and a chart
with its edges. In the tutorial: \href{https://seaborn.pydata.org/tutorial/axis_grids.html}{Building
structured multi-plot grids} (\code{FacetGrid}, \code{PairGrid}, \code{pairplot}, \code{JointGrid},
\code{jointplot}).

\textbf{In simplydata.} \code{draw\_one\_chart\_for\_each\_group\_in} with \code{down=} and
\code{across=} is \code{FacetGrid} with \code{row=} and \code{col=}; \code{same\_y=False} is
\code{sharey=False}. \code{sdata.draw\_all\_pairs} is \code{pairplot}. \code{draw\_side\_histograms}
is \code{jointplot}.

\begin{lstlisting}[style=py]
marks = sdata.get_example("marks")
chart = sdata.chart(marks)
chart.draw_dots(x="hours", y="marks")
chart.draw_one_chart_for_each_group_in(down="class", across="subject")
chart.save("grid.png")

pairs = sdata.draw_all_pairs(marks, colour="class")
pairs.save("pairs.png")
\end{lstlisting}

\textbf{Left to seaborn.} Your own function mapped on each panel of a grid (\code{FacetGrid.map}), and
different kinds of charts above and below the diagonal of a \code{PairGrid}.

\unsection{The look of a chart}{from-matplotlib:aesthetics}

\textbf{The job.} Choose the overall look, take away lines of the frame, and scale everything for a
paper, a talk or a poster. In the tutorial:
\href{https://seaborn.pydata.org/tutorial/aesthetics.html}{Controlling figure aesthetics}
(\code{set\_theme}, \code{set\_style}, \code{despine}, \code{set\_context}).

\textbf{In simplydata.} The normal look (\code{"paper"}) already leaves out the top and right lines,
as \code{despine} does. \code{use\_style} is \code{set\_style} and \code{set\_context} in one word:
\code{"presentation"} is like seaborn's \code{"talk"}. \code{hide\_frame} takes the whole frame away,
\code{show\_grid} and \code{hide\_grid} the grid lines, and \code{use\_background\_colour} the
background.

\begin{lstlisting}[style=py]
weather = sdata.get_example("weather")
chart = sdata.chart(weather)
chart.draw_lines(x="month", y="rain", colour="city")
chart.use_style("dark")
chart.show_grid()
chart.save("dark.png")
\end{lstlisting}

\unsection{Colour palettes}{from-matplotlib:palettes}

\textbf{The job.} Choose colours for groups (a qualitative palette), for numbers that go from low to
high (sequential) and for numbers on both sides of a middle (diverging). In the tutorial:
\href{https://seaborn.pydata.org/tutorial/color_palettes.html}{Choosing color palettes}.

\textbf{In simplydata.} The normal colours for groups are the Okabe-Ito set, safe for readers who do
not see some colours. \code{use\_colours} is \code{color\_palette} and \code{set\_palette}: a named set
(\code{"Set2"}, \code{"deep"}, \code{"colorblind"}), your own colours, or \code{how\_many=} colours
from a colour map. \code{use\_colour\_scale} is the \code{cmap} of a heat map; \code{"blue to red"}
with \code{middle=0} is a diverging palette centred on 0.

\begin{lstlisting}[style=py]
marks = sdata.get_example("marks")
chart = sdata.chart(marks)
chart.draw_bars(x="class", y="marks", colour="subject")
chart.use_colours("viridis", how_many=3)
chart.save("palette.png")
\end{lstlisting}

\unsection{The jobs left to matplotlib and seaborn}{from-matplotlib:left-out}

simplydata draws the charts most people need in research and teaching, with one word each. These jobs
are left out on purpose. Each one is a hand-over: \code{fig = chart.as\_matplotlib()}, then
matplotlib or seaborn code on \code{fig.axes}, then \code{fig.savefig(...)}.
\begin{itemize}
  \item \textbf{Arrows and vector fields} (\code{quiver}, \code{barbs}, \code{streamplot}).
  \item \textbf{Polygons, paths and other patches} (\code{fill}, \code{Polygon}), shown above.
  \item \textbf{A table drawn inside a chart} (\code{ax.table}).
  \item \textbf{Text of your own shape}: arrow styles, text along a path, many box styles (shown above).
  \item \textbf{Tick formatters as your own function} and scales of your own (shown above).
  \item \textbf{Axes placed by hand} (\code{fig.add\_axes}) and margins (\code{ax.margins}).
  \item \textbf{Images} (\code{imread}, \code{imsave}, \code{figimage}); images get their own words in
    a later phase.
  \item \textbf{Interaction}: backends, mouse and key events, widgets (\code{Slider}, \code{Button}).
  \item \textbf{More 3D}: voxels, arrows and contour sets in 3D.
  \item \textbf{seaborn's letter-value boxes} (\code{boxenplot}, shown above), a function of your own on
    each panel of a grid (\code{FacetGrid.map}) and the objects interface (\code{seaborn.objects}).
  \item \textbf{Any \code{rcParams} setting} beyond the looks, fonts and sizes of \code{use\_style},
    \code{use\_font} and \code{use\_font\_size} (shown above).
\end{itemize}

One more example: a vector field over a contour map.

\begin{lstlisting}[style=py]
island = sdata.get_example("island")
chart = sdata.chart(island)
chart.draw_contour(x="east", y="north", z="height")
fig = chart.as_matplotlib()
ax = fig.axes[0]
ax.quiver([2000, 4000], [1000, 3000], [1, -1], [1, 0.5], color="white")
fig.savefig("vectors.png")
\end{lstlisting}
).
% The topics follow the order of the matplotlib Quick start guide and User guide, then the seaborn tutorial, used as a
% checklist only: the text and the examples are our own, and each section links to the official page. Every example in
% a py box is run by tests/test_from_matplotlib_chapter.py (with "import simplydata as sdata" in front, and matplotlib's
% Agg backend), so it cannot go out of date. A line "% needs: plotly" just before a box says the box needs that package.

\unchapter{From matplotlib and seaborn to simplydata}{from-matplotlib}

This chapter is for you if you already know matplotlib or seaborn, or if you want to learn them after
simplydata. It walks through the topics of the official guides, in their order: the
\href{https://matplotlib.org/stable/users/explain/quick_start.html}{matplotlib Quick start guide} and
\href{https://matplotlib.org/stable/users/index.html}{User guide}, then the
\href{https://seaborn.pydata.org/tutorial.html}{seaborn tutorial}. For each topic it says what the job
is, which simplydata words do it, and a short example on the practice tables. Some jobs are left to
matplotlib and seaborn on purpose. For those, the chapter shows how to hand the chart over and keep
going in matplotlib.

Every example starts with \code{import simplydata as sdata}. It is written once here, not in each
box. The examples use the practice tables of \code{sdata.get\_example}, so you can run them as they
are.

\unsection{Five things to know first}{from-matplotlib:first}

\begin{enumerate}
  \item \textbf{A chart is a list of steps.} In matplotlib each call draws at once. In simplydata,
    \code{chart.draw\_dots(...)} and \code{chart.label\_x(...)} only write down a step. The chart is
    drawn when you save it, show it, or look at it in Jupyter. So the order of the steps rarely
    matters, and \code{chart.undo()} takes the last step back.
  \item \textbf{There is no figure and no axes to hold.} simplydata makes the matplotlib Figure and
    its Axes for you each time it draws. You never write \code{fig, ax = plt.subplots()}.
  \item \textbf{Columns by name.} You give the table once, to \code{sdata.chart(table)}, and then the
    column names, as seaborn does with \code{data=}. A word works on a pandas or a Polars table.
  \item \textbf{One word for one job.} matplotlib often has several ways to do the same thing
    (\code{plt.title}, \code{ax.set\_title}, \code{fig.suptitle}). simplydata has one word,
    \code{write\_title\_on\_top}, and \code{sdata.help("ax.set\_title")} names it.
  \item \textbf{The door is always open.} \code{fig = chart.as\_matplotlib()} gives the chart as a
    plain matplotlib Figure. Its charts are the Axes in \code{fig.axes}, so \code{ax = fig.axes[0]}.
    Then any matplotlib code works. It is a copy: changing \code{fig} does not change the chart.
    Save it with \code{fig.savefig(...)}.
\end{enumerate}

\begin{lstlisting}[style=py]
marks = sdata.get_example("marks")
chart = sdata.chart(marks)
chart.draw_dots(x="hours", y="marks")
fig = chart.as_matplotlib()              # a matplotlib Figure: use any matplotlib code here
ax = fig.axes[0]
ax.set_facecolor("whitesmoke")
fig.savefig("handed_over.png", dpi=150)
\end{lstlisting}

Do you know a matplotlib or seaborn name, but not the simplydata word? Ask:
\code{sdata.help("ax.set\_xlabel")}, \code{sdata.help("plt.hist")} or
\code{sdata.help("sns.lineplot")}. The name can have \code{plt.}, \code{ax.}, \code{fig.} or
\code{sns.} in front. The full list is in the section ``Coming from another tool?'' of the simplydata
dictionary
(\href{https://github.com/bioforkgithub/simplydata/blob/main/docs/manual/simplydataDictionary.pdf}%
{simplydataDictionary.pdf}).

\unsection{A first chart}{from-matplotlib:first-chart}

\textbf{The job.} Draw a line from two columns, give it a title and axis names, and save it as a
picture. In the guide: \href{https://matplotlib.org/stable/users/explain/quick_start.html}{Quick
start guide}, the simple example at its start.

\textbf{In simplydata.} \code{draw\_lines} is \code{ax.plot}; \code{write\_title\_on\_top},
\code{label\_x} and \code{label\_y} are \code{set\_title}, \code{set\_xlabel} and
\code{set\_ylabel}; \code{save} is \code{fig.savefig}. Without \code{label\_x}, the column name is the
axis name.

\begin{lstlisting}[style=py]
weather = sdata.get_example("weather")
chart = sdata.chart(weather)
chart.draw_lines(x="month", y="temperature", colour="city")
chart.write_title_on_top("Temperature through the year")
chart.label_y("Temperature (degrees C)")
chart.save("first_chart.png")
\end{lstlisting}

\unsection{The parts of a figure}{from-matplotlib:parts}

\textbf{The job.} Know the parts matplotlib draws: the Figure (the whole picture), the Axes (one
chart inside it), each Axis (the x and y scales) and the Artists (every line, dot, bar and piece of
text). In the guide: \href{https://matplotlib.org/stable/users/explain/quick_start.html}{Quick start
guide}, ``Parts of a Figure'', and
\href{https://matplotlib.org/stable/users/explain/artists/index.html}{Artists}.

\textbf{In simplydata.} A simplydata chart is one Figure with one Axes, or several Axes after
\code{draw\_one\_chart\_for\_each\_group\_in}. You do not touch the parts. You can read them back as
tables: \code{get\_numbers\_drawn} reads the lines, dots and bars, \code{get\_text\_drawn} the text,
and \code{get\_scales\_drawn} each axis. \code{as\_matplotlib} gives the parts themselves.

\begin{lstlisting}[style=py]
marks = sdata.get_example("marks")
chart = sdata.chart(marks)
chart.draw_bars(x="class", y="marks")
chart.get_numbers_drawn()                # one row for each bar: where it is and how high
chart.get_scales_drawn()                 # the range and kind of each axis
fig = chart.as_matplotlib()
print(len(fig.axes), "Axes in the Figure")
\end{lstlisting}

\unsection{The data a chart takes}{from-matplotlib:data}

\textbf{The job.} Give the numbers to draw. matplotlib takes lists or NumPy arrays, or column names
with \code{data=}. In the guide: \href{https://matplotlib.org/stable/users/explain/quick_start.html}%
{Quick start guide}, ``Types of inputs to plotting functions''.

\textbf{In simplydata.} A table and column names is the usual way: \code{sdata.chart(table)} reads a
table, a file name or a pandas or Polars DataFrame. Without a table, give lists to the step itself.
Numbers made with \code{sdata.number\_list} start their own chart.

\begin{lstlisting}[style=py]
chart = sdata.chart()                    # no table: lists go to the step
chart.draw_lines(x=[1, 2, 3, 4], y=[10, 12, 9, 14])
chart.save("from_lists.png")

numbers = sdata.number_list([4, 8, 15, 16, 23, 42])
chart = numbers.start_chart()
chart.draw_histogram(x="value")
chart.save("from_numbers.png")
\end{lstlisting}

\unsection{Two ways of writing, one in simplydata}{from-matplotlib:styles}

\textbf{The job.} matplotlib can be written in two styles: the explicit one (\code{fig, ax =
plt.subplots()}, then \code{ax.plot(...)}) and the implicit one with \code{pyplot}
(\code{plt.plot(...)}, which draws on the current Axes). In the guide:
\href{https://matplotlib.org/stable/users/explain/quick_start.html}{Quick start guide}, ``Coding
styles''.

\textbf{In simplydata.} There is one style. The chart is the object you name, as in the explicit style,
and nothing is ``current''. So two charts never draw on each other. To make many charts the same way,
write a small function that gives back a chart, as the guide suggests for matplotlib.

\begin{lstlisting}[style=py]
def marks_chart(subject):
    marks = sdata.get_example("marks")
    marks.keep_rows_where("subject", is_equal_to=subject)
    chart = sdata.chart(marks)
    chart.draw_dots(x="hours", y="marks")
    chart.write_title_on_top(subject)
    return chart


marks_chart("Maths").save("maths.png")
marks_chart("Science").save("science.png")
\end{lstlisting}

\unsection{Colours, lines and dots}{from-matplotlib:styling}

\textbf{The job.} Choose the colour, width and dash of a line, and the size, shape and transparency of
dots. In the guide: \href{https://matplotlib.org/stable/users/explain/quick_start.html}{Quick start
guide}, ``Styling Artists'', and
\href{https://matplotlib.org/stable/users/explain/colors/colors.html}{Specifying colors}.

\textbf{In simplydata.} Options with plain names on the draw words: \code{thickness=} is
\code{linewidth}, \code{dashed=True} and \code{dotted=True} are \code{linestyle}, \code{size=} is the
dot size (a column, for sizes that change), \code{transparency=} is one minus \code{alpha}.
\code{use\_colour} sets one colour, \code{use\_colours} one colour for each group, and
\code{use\_shapes} a shape for each group. The groups come from \code{colour="column"}, as seaborn's
\code{hue=}.

\begin{lstlisting}[style=py]
weather = sdata.get_example("weather")
chart = sdata.chart(weather)
chart.draw_lines(x="month", y="temperature", colour="city", thickness=3, dashed="city")
chart.use_colours("black", "darkorange")
chart.save("styled_lines.png")

marks = sdata.get_example("marks")
chart = sdata.chart(marks)
chart.draw_dots(x="hours", y="marks", colour="class", size="height", transparency=0.3)
chart.use_shapes("circle", "square", "triangle")
chart.save("styled_dots.png")
\end{lstlisting}

\unsection{Titles, labels, notes and legends}{from-matplotlib:labels}

\textbf{The job.} Name the axes, write text on the chart, point at a place with an arrow, and place
the legend. In the guide: \href{https://matplotlib.org/stable/users/explain/quick_start.html}{Quick
start guide}, ``Labelling plots'';
\href{https://matplotlib.org/stable/users/explain/text/text_intro.html}{Text in Matplotlib},
\href{https://matplotlib.org/stable/users/explain/text/annotations.html}{Annotations} and the
\href{https://matplotlib.org/stable/users/explain/axes/legend_guide.html}{Legend guide}.

\textbf{In simplydata.} \code{label\_x}, \code{label\_y}, \code{write\_title\_on\_top} and
\code{write\_caption\_at\_bottom} write the text around the chart. \code{write\_note} is
\code{ax.text}, and with \code{with\_arrow=True} it is \code{ax.annotate}. The text goes to
matplotlib as you write it, so matplotlib's maths between dollar signs works too. A legend appears by
itself when \code{colour=} makes groups; \code{put\_legend\_inside}, \code{put\_legend\_at\_bottom}
and the other \code{put\_legend\_} words place it, \code{label\_legend} names it,
\code{order\_legend} orders it and \code{hide\_legend} hides it.

\begin{lstlisting}[style=py]
growth = sdata.get_example("growth")
chart = sdata.chart(growth)
chart.draw_lines(x="year", y="population", colour="country")
chart.label_y("People (millions)")
chart.label_x("Year, $t$")
chart.write_note("fastest growth", x=2020, y=12, with_arrow=True)
chart.write_caption_at_bottom("Made-up data")
chart.put_legend_inside("top left")
chart.label_legend("Country")
chart.save("labels.png")
\end{lstlisting}

\textbf{Left to matplotlib.} Arrows with your own shape and connection, text along a path, and
boxes of many styles. Hand the chart over:

\begin{lstlisting}[style=py]
growth = sdata.get_example("growth")
chart = sdata.chart(growth)
chart.draw_lines(x="year", y="population", colour="country")
fig = chart.as_matplotlib()
ax = fig.axes[0]
ax.annotate("a bend", xy=(2010, 6), xytext=(2003, 10),
            arrowprops={"arrowstyle": "-|>", "connectionstyle": "arc3,rad=0.3"},
            bbox={"boxstyle": "round", "facecolor": "lightyellow"})
fig.savefig("own_arrow.png")
\end{lstlisting}

\unsection{Scales, ticks, dates and groups on an axis}{from-matplotlib:axes}

\textbf{The job.} Choose the range and the scale of an axis (log, square root), where the ticks go and
how their numbers are written, dates and text on an axis, and a second axis. In the guide:
\href{https://matplotlib.org/stable/users/explain/quick_start.html}{Quick start guide}, ``Axis
scales and ticks'';
\href{https://matplotlib.org/stable/users/explain/axes/axes_scales.html}{Axis scales} and
\href{https://matplotlib.org/stable/users/explain/axes/axes_ticks.html}{Axis ticks}.

\textbf{In simplydata.} \code{set\_range\_x} and \code{set\_range\_y} are \code{set\_xlim} and
\code{set\_ylim}; \code{use\_log\_scale\_y} and \code{use\_scale\_y("square root")} are
\code{set\_yscale}; \code{put\_ticks\_x} is \code{set\_xticks} (with \code{labels=} it is
\code{set\_xticklabels}, and \code{every=} a \code{MultipleLocator}); \code{write\_ticks\_y\_as}
chooses a formatter; \code{rotate\_labels\_x} turns the labels; \code{flip\_y} is
\code{invert\_yaxis}. Dates and text on an axis need nothing: simplydata finds them. Months and
weekdays come in calendar order, and \code{order\_x} sets any other order. \code{add\_second\_y} is
\code{twinx}, and \code{add\_second\_y(formula=...)} is \code{secondary\_yaxis}.

\begin{lstlisting}[style=py]
shop = sdata.get_example("shop")
chart = sdata.chart(shop)
chart.draw_lines(x="date", y="sales", colour="shop")
chart.write_ticks_x_as("Jan 2026")
chart.write_ticks_y_as("thousands")
chart.rotate_labels_x(45)
chart.save("dates.png")

marks = sdata.get_example("marks")
chart = sdata.chart(marks)
chart.draw_bars(x="class", y="marks")
chart.order_x("C", "B", "A")
chart.set_range_y(0, 100)
chart.put_ticks_y(0, 25, 50, 75, 100)
chart.save("groups.png")

weather = sdata.get_example("weather")
chart = sdata.chart(weather)
chart.draw_lines(x="month", y="temperature", colour="city")
chart.add_second_y(formula="y * 9 / 5 + 32")
chart.save("second_axis.png")
\end{lstlisting}

\textbf{Left to matplotlib.} A tick formatter written as your own function, and scales of your own
(\code{FuncScale}). Hand the chart over:

\begin{lstlisting}[style=py]
from matplotlib.ticker import FuncFormatter

marks = sdata.get_example("marks")
chart = sdata.chart(marks)
chart.draw_dots(x="hours", y="marks")
fig = chart.as_matplotlib()
fig.axes[0].xaxis.set_major_formatter(FuncFormatter(lambda value, place: f"{value:.0f} h"))
fig.savefig("own_ticks.png")
\end{lstlisting}

\unsection{Numbers as colours}{from-matplotlib:colour-maps}

\textbf{The job.} Show a number as a colour (a heat map or a contour map), choose the colour map and
how numbers map to colours, and add a colour bar. In the guide:
\href{https://matplotlib.org/stable/users/explain/quick_start.html}{Quick start guide}, ``Color
mapped data''; \href{https://matplotlib.org/stable/users/explain/colors/colormaps.html}{Choosing
Colormaps} and \href{https://matplotlib.org/stable/users/explain/colors/colormapnorms.html}{Colormap
normalization}.

\textbf{In simplydata.} \code{draw\_heatmap} is \code{imshow} or \code{pcolormesh} (seaborn
\code{heatmap}), and \code{draw\_contour} is \code{contourf}, also for scattered points.
\code{use\_colour\_scale} is \code{cmap=}: it chooses the colour map. Its options are the
normalizations: \code{middle=0} puts the middle colour at 0 (a \code{TwoSlopeNorm}), \code{log=True}
spreads the colours on a log scale, \code{steps=} cuts the scale into bands, and \code{between=} sets
\code{vmin} and \code{vmax}. The colour bar comes by itself; \code{label\_colour\_bar} names it and
\code{hide\_colour\_bar} hides it.

\begin{lstlisting}[style=py]
weather = sdata.get_example("weather")
chart = sdata.chart(weather)
chart.draw_heatmap(x="month", y="city", value="rain", with_numbers=True)
chart.use_colour_scale("blues", between=[0, 300])
chart.label_colour_bar("Rain (mm)")
chart.save("heatmap.png")

island = sdata.get_example("island")
chart = sdata.chart(island)
chart.draw_contour(x="east", y="north", z="height")
chart.use_colour_scale("terrain", steps=8)
chart.save("contour.png")
\end{lstlisting}

\unsection{Several charts in one picture}{from-matplotlib:several}

\textbf{The job.} Put several Axes in one Figure: a grid of charts, one chart for each group, charts of
different sizes, or a small chart inside a big one. In the guide:
\href{https://matplotlib.org/stable/users/explain/quick_start.html}{Quick start guide}, ``Working
with multiple Figures and Axes'';
\href{https://matplotlib.org/stable/users/explain/axes/arranging_axes.html}{Arranging multiple Axes in
a Figure} and the
\href{https://matplotlib.org/stable/users/explain/axes/constrainedlayout_guide.html}{Constrained layout
guide}.

\textbf{In simplydata.} \code{draw\_one\_chart\_for\_each\_group\_in} is seaborn's
\code{FacetGrid} (\code{col=}, \code{row=}, \code{col\_wrap=}). \code{sdata.put\_side\_by\_side},
\code{sdata.put\_one\_above\_other} and \code{sdata.put\_in\_rows\_and\_columns} are
\code{plt.subplots} and \code{subplot\_mosaic} for charts you made one by one; a nested list makes
a chart span several places, and \code{with\_letters=True} writes A, B, C. \code{put\_chart\_inside} is
\code{inset\_axes}, and \code{add\_zoom\_box} is \code{indicate\_inset\_zoom}. The layout is always
matplotlib's constrained layout, so labels do not overlap.

\begin{lstlisting}[style=py]
marks = sdata.get_example("marks")
chart = sdata.chart(marks)
chart.draw_dots(x="hours", y="marks")
chart.draw_one_chart_for_each_group_in("subject", per_row=3)
chart.save("one_per_group.png")

dots = sdata.chart(marks).draw_dots(x="hours", y="marks")
bars = sdata.chart(marks).draw_bars(x="class", y="marks")
boxes = sdata.chart(marks).draw_boxes(x="subject", y="marks")
page = sdata.put_in_rows_and_columns([[dots, dots], [bars, boxes]], with_letters=True)
page.save("page.png")
\end{lstlisting}

\textbf{Left to matplotlib.} Axes placed by hand at any place in the Figure (\code{fig.add\_axes}).
\code{page.as\_matplotlib()} gives the page as one Figure: each chart is a SubFigure in
\code{fig.subfigs}, and \code{fig.axes} holds every Axes of every chart.

\unsection{Looks and settings}{from-matplotlib:style}

\textbf{The job.} Change the look of every chart at once: a style sheet, the fonts and the sizes of the
text (matplotlib's \code{rcParams}). In the guide:
\href{https://matplotlib.org/stable/users/explain/customizing.html}{Customizing Matplotlib with style
sheets and rcParams}.

\textbf{In simplydata.} \code{use\_style} picks one of five looks: \code{"paper"} (normal),
\code{"school"}, \code{"presentation"}, \code{"dark"} and \code{"black and white"}.
\code{use\_font} and \code{use\_font\_size} set the text. Each chart keeps its own look; nothing
changes matplotlib's settings for other code.

\begin{lstlisting}[style=py]
marks = sdata.get_example("marks")
chart = sdata.chart(marks)
chart.draw_boxes(x="class", y="marks")
chart.use_style("presentation")
chart.use_font("serif")
chart.use_font_size(16, title=22)
chart.save("presentation.png")
\end{lstlisting}

\textbf{Left to matplotlib.} Any other setting of \code{rcParams}. A setting such as the colour of the
axis lines is used while matplotlib draws, so change the drawn parts after the hand-over:

\begin{lstlisting}[style=py]
marks = sdata.get_example("marks")
chart = sdata.chart(marks)
chart.draw_boxes(x="class", y="marks")
fig = chart.as_matplotlib()
for side in fig.axes[0].spines.values():
    side.set_color("steelblue")
fig.savefig("own_settings.png")
\end{lstlisting}

\unsection{Lines, bands and shapes on a chart}{from-matplotlib:marks}

\textbf{The job.} Draw a reference line, shade a range, draw a rectangle, a circle or an arrow. In the
guide: \href{https://matplotlib.org/stable/plot_types/index.html}{Plot types} and
\href{https://matplotlib.org/stable/users/explain/text/annotations.html}{Annotations}.

\textbf{In simplydata.} \code{draw\_line\_at\_y} is \code{axhline} (with \code{between=} it is
\code{hlines}), \code{draw\_line\_at\_x} is \code{axvline}, \code{draw\_line\_with\_slope} is
\code{axline}, \code{shade\_between\_x} is \code{axvspan}, \code{draw\_band} is
\code{fill\_between}, \code{draw\_rectangle} and \code{draw\_circle} are the \code{Rectangle} and
\code{Circle} patches, and \code{draw\_arrow} is an arrow from one point to another.

\begin{lstlisting}[style=py]
marks = sdata.get_example("marks")
chart = sdata.chart(marks)
chart.draw_dots(x="hours", y="marks")
chart.draw_line_at_y(50, labelled="pass mark")
chart.shade_between_x(6, 8)
chart.draw_rectangle(x_between=[3, 5], y_between=[60, 80])
chart.draw_arrow(start=[2, 40], end=[4, 60])
chart.save("shapes.png")
\end{lstlisting}

\textbf{Left to matplotlib.} Polygons, paths and other patches. Hand the chart over and add them:

\begin{lstlisting}[style=py]
from matplotlib.patches import Polygon

marks = sdata.get_example("marks")
chart = sdata.chart(marks)
chart.draw_dots(x="hours", y="marks")
fig = chart.as_matplotlib()
fig.axes[0].add_patch(Polygon([[2, 40], [5, 90], [8, 60]], alpha=0.2, color="purple"))
fig.savefig("polygon.png")
\end{lstlisting}

\unsection{Charts in 3D}{from-matplotlib:3d}

\textbf{The job.} Draw a surface, dots or lines in three dimensions and choose the view. In the guide:
\href{https://matplotlib.org/stable/users/explain/toolkits/mplot3d.html}{The mplot3d toolkit}.

\textbf{In simplydata.} \code{draw\_surface} is \code{plot\_surface} (with \code{as\_wires=True} it is
\code{plot\_wireframe}), \code{draw\_dots\_3d} is \code{scatter} on 3D Axes, \code{draw\_lines\_3d}
is \code{plot} in 3D and \code{draw\_bars\_3d} is \code{bar3d}. \code{look\_from} is
\code{view\_init}, and \code{label\_z}, \code{set\_range\_z} and \code{put\_ticks\_z} work on the third
axis.

\begin{lstlisting}[style=py]
island = sdata.get_example("island")
chart = sdata.chart(island)
chart.draw_surface(x="east", y="north", z="height")
chart.look_from(height=40, around=30)
chart.label_z("Height (m)")
chart.save("surface.png")
\end{lstlisting}

\unsection{Moving charts}{from-matplotlib:animations}

\textbf{The job.} Make a chart that moves, one picture after another. In the guide:
\href{https://matplotlib.org/stable/users/explain/animations/animations.html}{Animations using
Matplotlib}.

\textbf{In simplydata.} \code{animate(over="year")} draws one picture for each value of a column:
matplotlib draws each picture and Pillow writes them into a .gif file.

\begin{lstlisting}[style=py]
growth = sdata.get_example("growth")
chart = sdata.chart(growth)
chart.draw_bars(x="country", y="population")
chart.animate(over="year")
chart.save("growth.gif")
\end{lstlisting}

\unsection{Saving, showing and web pages}{from-matplotlib:saving}

\textbf{The job.} Save the picture in a file type and size you need, show it in a window, or make a
chart you can point at in a web page. In the guide:
\href{https://matplotlib.org/stable/users/explain/figure/index.html}{Figures and backends}, with
\href{https://matplotlib.org/stable/users/explain/figure/backends.html}{Backends} and
\href{https://matplotlib.org/stable/users/explain/figure/interactive.html}{Interactive figures}.

\textbf{In simplydata.} \code{save} is \code{savefig}: the end of the file name picks the type
(.png, .pdf, .svg, .eps, .tif, .jpg, .gif or .html), the picture is 300 dots per inch and trimmed to
what is drawn, \code{resolution=} is \code{dpi}, and \code{see\_through=True} is
\code{transparent=True}. \code{set\_size} is \code{figsize}, in inches or cm. \code{show} is
\code{plt.show}. A .html file and \code{show\_in\_web\_page} make a web page, drawn by plotly from
what matplotlib drew; \code{show\_when\_pointing} chooses what it shows when the pointer rests on a
dot. In Jupyter a chart shows by itself.

\begin{lstlisting}[style=py]
marks = sdata.get_example("marks")
chart = sdata.chart(marks)
chart.draw_dots(x="hours", y="marks", colour="class")
chart.set_size(width=17, height=10, unit="cm")
chart.save("for_the_journal.pdf")
chart.save("for_a_slide.png", resolution=150, see_through=True)
chart.get_size(unit="cm")
\end{lstlisting}

% needs: plotly
\begin{lstlisting}[style=py]
marks = sdata.get_example("marks")
chart = sdata.chart(marks)
chart.draw_dots(x="hours", y="marks", colour="class")
chart.show_when_pointing("name", "subject")
chart.save("marks.html")
\end{lstlisting}

\textbf{The hand-over and web pages.} \code{chart.as\_matplotlib()} always gives the matplotlib
Figure, also for a chart you save as a web page: the web page is a copy that plotly draws from it.
Changes you make in matplotlib after the hand-over are not in \code{chart.save("marks.html")}.

\textbf{Left to matplotlib.} Choosing the backend, windows that react to the mouse or the keyboard
(events), and widgets such as sliders and buttons.

\unsection{Relationships: dots and lines}{from-matplotlib:relational}

\textbf{The job.} Show how one number goes with another: dots for each row, or a line through the
average at each x, with groups by colour, size or shape. In the tutorial:
\href{https://seaborn.pydata.org/tutorial/relational.html}{Visualizing statistical relationships}
(\code{relplot}, \code{scatterplot}, \code{lineplot}), and the overview of the
\href{https://seaborn.pydata.org/tutorial/function_overview.html}{plotting functions} and the
\href{https://seaborn.pydata.org/tutorial/data_structure.html}{data structures} seaborn takes.

\textbf{In simplydata.} \code{draw\_dots} is \code{scatterplot}: \code{colour=} is \code{hue=},
\code{size=} is \code{size=} and \code{shape=} is \code{style=}. \code{draw\_lines} is
\code{lineplot}; it draws the average of the rows at each x, and \code{error=} adds the band.
\code{relplot}'s panels are \code{draw\_one\_chart\_for\_each\_group\_in}.

\begin{lstlisting}[style=py]
marks = sdata.get_example("marks")
chart = sdata.chart(marks)
chart.draw_dots(x="hours", y="marks", colour="class", shape="subject")
chart.save("relationship.png")

school = sdata.get_example("school")
chart = sdata.chart(school)
chart.draw_lines(x="sleep", y="marks", colour="breakfast", error="confidence interval")
chart.save("average_line.png")
\end{lstlisting}

\unsection{Distributions}{from-matplotlib:distributions}

\textbf{The job.} Show how the values of a column are spread: a histogram, a smooth density, the
share of values up to each point, a rug, and two columns together. In the tutorial:
\href{https://seaborn.pydata.org/tutorial/distributions.html}{Visualizing distributions of data}
(\code{displot}, \code{histplot}, \code{kdeplot}, \code{ecdfplot}, \code{rugplot}).

\textbf{In simplydata.} \code{draw\_histogram} is \code{histplot} (\code{with\_bins=},
\code{show="density"}, \code{stacked=True}, \code{cumulative=True}); \code{draw\_density} is
\code{kdeplot}; \code{draw\_cumulative\_share} is \code{ecdfplot}; \code{draw\_rug\_lines} is
\code{rugplot}. With \code{y=} as well, \code{draw\_histogram} and \code{draw\_density} show two
columns together, and \code{draw\_side\_histograms} adds the edges of \code{jointplot}.
\code{fit\_distribution} draws a fitted curve over a histogram.

\begin{lstlisting}[style=py]
marks = sdata.get_example("marks")
chart = sdata.chart(marks)
chart.draw_histogram(x="marks", colour="class", show="density")
chart.draw_density(x="marks", colour="class")
chart.draw_rug_lines(x="marks")
chart.save("spread.png")

chart = sdata.chart(marks)
chart.draw_dots(x="hours", y="marks")
chart.draw_side_histograms(drawn_as="density")
chart.save("joint.png")
\end{lstlisting}

\unsection{Groups: categorical charts}{from-matplotlib:categorical}

\textbf{The job.} Compare the values of groups: every dot, boxes, violins, bars of an average with
an interval, points joined across groups, and counts. In the tutorial:
\href{https://seaborn.pydata.org/tutorial/categorical.html}{Visualizing categorical data}
(\code{catplot}, \code{stripplot}, \code{swarmplot}, \code{boxplot}, \code{violinplot},
\code{barplot}, \code{pointplot}, \code{countplot}).

\textbf{In simplydata.} \code{draw\_all\_dots} is \code{stripplot} and \code{swarmplot};
\code{draw\_boxes} is \code{boxplot}; \code{draw\_violins} is \code{violinplot} (\code{in\_halves=}
is \code{split=True}); \code{draw\_bars} is \code{barplot}, with \code{show=} for the average,
median, total or count; \code{draw\_interaction\_lines} is \code{pointplot}; \code{draw\_count\_bars}
is \code{countplot}. \code{catplot}'s panels are \code{draw\_one\_chart\_for\_each\_group\_in}.

\begin{lstlisting}[style=py]
marks = sdata.get_example("marks")
chart = sdata.chart(marks)
chart.draw_boxes(x="subject", y="marks")
chart.draw_all_dots(x="subject", y="marks")
chart.save("boxes_and_dots.png")

school = sdata.get_example("school")
chart = sdata.chart(school)
chart.draw_interaction_lines(x="class", y="marks", colour="breakfast")
chart.save("points.png")

chart = sdata.chart(school)
chart.draw_count_bars(x="class", colour="breakfast")
chart.save("counts.png")
\end{lstlisting}

\textbf{Left to seaborn.} \code{boxenplot}, the letter-value boxes for big tables. Hand an empty chart
over and let seaborn draw on its Axes:

\begin{lstlisting}[style=py]
import seaborn as sns

marks = sdata.get_example("marks")
chart = sdata.chart(marks)
chart.label_x("subject")
fig = chart.as_matplotlib()
sns.boxenplot(marks.as_pandas(), x="subject", y="marks", ax=fig.axes[0])
fig.savefig("boxen.png")
\end{lstlisting}

\unsection{Averages with error bars}{from-matplotlib:error-bars}

\textbf{The job.} Show how sure an average is (a confidence interval or a standard error) or how
spread the values are (a standard deviation). In the tutorial:
\href{https://seaborn.pydata.org/tutorial/error_bars.html}{Statistical estimation and error bars}.

\textbf{In simplydata.} \code{error=} on \code{draw\_bars}, \code{draw\_lines} and \code{draw\_dots}
takes \code{"confidence interval"}, \code{"standard error"} or \code{"standard deviation"}, as
seaborn's \code{errorbar=("ci", 95)}, \code{"se"} and \code{"sd"}. \code{draw\_error\_bars} adds them
to dots or lines. The interval is the t interval, not seaborn's bootstrap, so the same table always
gives the same bars, and a line under the chart says which one was drawn.

\begin{lstlisting}[style=py]
marks = sdata.get_example("marks")
chart = sdata.chart(marks)
chart.draw_bars(x="class", y="marks", show="average", error="confidence interval")
chart.save("error_bars.png")
\end{lstlisting}

\unsection{Lines that fit}{from-matplotlib:regression}

\textbf{The job.} Draw the line that fits the dots best, with its confidence band, a curve, a smooth
line, and what is left over. In the tutorial:
\href{https://seaborn.pydata.org/tutorial/regression.html}{Estimating regression fits}
(\code{regplot}, \code{lmplot}, \code{residplot}).

\textbf{In simplydata.} \code{fit\_line} is \code{regplot}; \code{show\_equation=True} writes the
equation and $R^2$ on the chart. \code{fit\_curve(of\_degree=2)} is \code{order=2}.
\code{draw\_smooth\_line} is \code{lowess=True}. \code{draw\_residual\_dots} is \code{residplot}.
For \code{lmplot}'s panels, add \code{draw\_one\_chart\_for\_each\_group\_in}.

\begin{lstlisting}[style=py]
marks = sdata.get_example("marks")
chart = sdata.chart(marks)
chart.draw_dots(x="hours", y="marks")
chart.fit_line(show_equation=True)
chart.save("fitted_line.png")

chart = sdata.chart(marks)
chart.draw_residual_dots(x="hours", y="marks")
chart.save("left_over.png")
\end{lstlisting}

\unsection{Grids of charts}{from-matplotlib:grids}

\textbf{The job.} One chart for each group in a grid, every column against every other, and a chart
with its edges. In the tutorial: \href{https://seaborn.pydata.org/tutorial/axis_grids.html}{Building
structured multi-plot grids} (\code{FacetGrid}, \code{PairGrid}, \code{pairplot}, \code{JointGrid},
\code{jointplot}).

\textbf{In simplydata.} \code{draw\_one\_chart\_for\_each\_group\_in} with \code{down=} and
\code{across=} is \code{FacetGrid} with \code{row=} and \code{col=}; \code{same\_y=False} is
\code{sharey=False}. \code{sdata.draw\_all\_pairs} is \code{pairplot}. \code{draw\_side\_histograms}
is \code{jointplot}.

\begin{lstlisting}[style=py]
marks = sdata.get_example("marks")
chart = sdata.chart(marks)
chart.draw_dots(x="hours", y="marks")
chart.draw_one_chart_for_each_group_in(down="class", across="subject")
chart.save("grid.png")

pairs = sdata.draw_all_pairs(marks, colour="class")
pairs.save("pairs.png")
\end{lstlisting}

\textbf{Left to seaborn.} Your own function mapped on each panel of a grid (\code{FacetGrid.map}), and
different kinds of charts above and below the diagonal of a \code{PairGrid}.

\unsection{The look of a chart}{from-matplotlib:aesthetics}

\textbf{The job.} Choose the overall look, take away lines of the frame, and scale everything for a
paper, a talk or a poster. In the tutorial:
\href{https://seaborn.pydata.org/tutorial/aesthetics.html}{Controlling figure aesthetics}
(\code{set\_theme}, \code{set\_style}, \code{despine}, \code{set\_context}).

\textbf{In simplydata.} The normal look (\code{"paper"}) already leaves out the top and right lines,
as \code{despine} does. \code{use\_style} is \code{set\_style} and \code{set\_context} in one word:
\code{"presentation"} is like seaborn's \code{"talk"}. \code{hide\_frame} takes the whole frame away,
\code{show\_grid} and \code{hide\_grid} the grid lines, and \code{use\_background\_colour} the
background.

\begin{lstlisting}[style=py]
weather = sdata.get_example("weather")
chart = sdata.chart(weather)
chart.draw_lines(x="month", y="rain", colour="city")
chart.use_style("dark")
chart.show_grid()
chart.save("dark.png")
\end{lstlisting}

\unsection{Colour palettes}{from-matplotlib:palettes}

\textbf{The job.} Choose colours for groups (a qualitative palette), for numbers that go from low to
high (sequential) and for numbers on both sides of a middle (diverging). In the tutorial:
\href{https://seaborn.pydata.org/tutorial/color_palettes.html}{Choosing color palettes}.

\textbf{In simplydata.} The normal colours for groups are the Okabe-Ito set, safe for readers who do
not see some colours. \code{use\_colours} is \code{color\_palette} and \code{set\_palette}: a named set
(\code{"Set2"}, \code{"deep"}, \code{"colorblind"}), your own colours, or \code{how\_many=} colours
from a colour map. \code{use\_colour\_scale} is the \code{cmap} of a heat map; \code{"blue to red"}
with \code{middle=0} is a diverging palette centred on 0.

\begin{lstlisting}[style=py]
marks = sdata.get_example("marks")
chart = sdata.chart(marks)
chart.draw_bars(x="class", y="marks", colour="subject")
chart.use_colours("viridis", how_many=3)
chart.save("palette.png")
\end{lstlisting}

\unsection{The jobs left to matplotlib and seaborn}{from-matplotlib:left-out}

simplydata draws the charts most people need in research and teaching, with one word each. These jobs
are left out on purpose. Each one is a hand-over: \code{fig = chart.as\_matplotlib()}, then
matplotlib or seaborn code on \code{fig.axes}, then \code{fig.savefig(...)}.
\begin{itemize}
  \item \textbf{Arrows and vector fields} (\code{quiver}, \code{barbs}, \code{streamplot}).
  \item \textbf{Polygons, paths and other patches} (\code{fill}, \code{Polygon}), shown above.
  \item \textbf{A table drawn inside a chart} (\code{ax.table}).
  \item \textbf{Text of your own shape}: arrow styles, text along a path, many box styles (shown above).
  \item \textbf{Tick formatters as your own function} and scales of your own (shown above).
  \item \textbf{Axes placed by hand} (\code{fig.add\_axes}) and margins (\code{ax.margins}).
  \item \textbf{Images} (\code{imread}, \code{imsave}, \code{figimage}); images get their own words in
    a later phase.
  \item \textbf{Interaction}: backends, mouse and key events, widgets (\code{Slider}, \code{Button}).
  \item \textbf{More 3D}: voxels, arrows and contour sets in 3D.
  \item \textbf{seaborn's letter-value boxes} (\code{boxenplot}, shown above), a function of your own on
    each panel of a grid (\code{FacetGrid.map}) and the objects interface (\code{seaborn.objects}).
  \item \textbf{Any \code{rcParams} setting} beyond the looks, fonts and sizes of \code{use\_style},
    \code{use\_font} and \code{use\_font\_size} (shown above).
\end{itemize}

One more example: a vector field over a contour map.

\begin{lstlisting}[style=py]
island = sdata.get_example("island")
chart = sdata.chart(island)
chart.draw_contour(x="east", y="north", z="height")
fig = chart.as_matplotlib()
ax = fig.axes[0]
ax.quiver([2000, 4000], [1000, 3000], [1, -1], [1, 0.5], color="white")
fig.savefig("vectors.png")
\end{lstlisting}
